\chapter{The Grassmannian over $\mathbb{Z}$: affine charts and gluing}
\label{chap:Picard_GrassmannianCells}

\section{Set-up: charts indexed by $d$-subsets}
\label{sec:gr_setup}

Fix integers \(r \ge d \ge 1\). For a \(d \times r\) matrix \(M\) over a ring
and a subset \(I \subseteq \{1, \dots, r\}\) with \(\#(I) = d\), write \(M_I\)
for the \emph{\(I\)-th minor} of \(M\): the \(d \times d\) submatrix whose
columns are those indexed by \(I\). We construct the Grassmannian scheme
\(\mathrm{Gr}(r,d)\) over \(\mathbb{Z}\) by gluing one affine chart \(U^I\) for
each such \(I\). The construction is, conceptually, the quotient
\(\mathrm{GL}_{d,\mathbb{Z}} \backslash V\) of the scheme \(V\) of rank-\(d\)
matrices, but we never use the language of group-scheme actions: the charts and
their transition maps are elementary matrix algebra over \(\mathbb{Z}\). Taking
\(d = 1\) recovers \(\mathrm{Gr}(r,1) = \mathbb{P}^{r-1}_{\mathbb{Z}}\).

\begin{definition}[Affine chart \(U^I\)]\leanok
  \label{def:gr_affine_chart}
  \dcref{custom}
  \lean{AlgebraicGeometry.Grassmannian.affineChart}
  % SOURCE: \cite{nitsure-hilbert-quot}, "Construction of Hilbert and Quot Schemes", \S 1,
  % sub-section "Construction by gluing together affine patches" (read from
  % references/nitsure-hilbert-quot-src/nitsure-hilbert-quot.tex, L807-L821).
  % SOURCE QUOTE:
  % "If $M$ is a $d\times r$-matrix, and
  %  $I\subset \{1,\ldots, r\}$ with cardinality $\#(I)$ equal
  %  to $d$, the $I$ th minor $M_I$ of $M$ will
  %  mean the $d\times d$ minor of $M$ whose columns are indexed by
  %  $I$.
  %
  %  For any subset $I\subset \{1,\ldots, r\}$ with $\#(I) =d$,
  %  consider the $d\times r$ matrix $X^I$ whose $I$ the minor
  %  $X^I_I$ is the $d\times d$ identity matrix $1_{d\times d}$, while the remaining
  %  entries of $X^I$ are independent variables $x^I_{p,q}$ over $\Z$.
  %  Let $\Z[X^I]$ denote the polynomial ring in the variables
  %  $x^I_{p,q}$, and let $U^I = \spec \Z[X^I]$, which is non-canonically
  %  isomorphic to the affine space $\AA^{d(r-d)}_{\Z}$."
  \textit{Source: \cite{nitsure-hilbert-quot}, \S 1.}
  Let \(I \subseteq \{1, \dots, r\}\) with \(\#(I) = d\). Let \(X^I\) denote the
  \(d \times r\) matrix whose \(I\)-th minor \(X^I_I\) is the \(d \times d\)
  identity matrix \(1_{d \times d}\) and whose remaining \(d(r - d)\) entries are
  independent indeterminates \(x^I_{p,q}\) (\(1 \le p \le d\),
  \(q \in \{1,\dots,r\} \setminus I\)) over \(\mathbb{Z}\). Write
  \(\mathbb{Z}[X^I]\) for the polynomial ring in these \(d(r-d)\) variables and
  set
  \[
    U^I \;:=\; \Spec \mathbb{Z}[X^I].
  \]
  The choice of an ordering of the free entries gives a (non-canonical)
  isomorphism \(U^I \cong \mathbb{A}^{d(r-d)}_{\mathbb{Z}}\); in particular
  \(U^I\) is a smooth affine \(\mathbb{Z}\)-scheme of relative dimension
  \(d(r-d)\). Intuitively a \(\mathbb{Z}\)-point of \(U^I\) is a \(d \times r\)
  matrix whose \(I\)-columns form the identity, i.e.\ a rank-\(d\) quotient of
  \(\mathbb{Z}^{r}\) trivialised on the coordinates \(I\).
\end{definition}

\section{Transition maps and the cocycle condition}
\label{sec:gr_transition}

\subsection*{Mathlib inputs for the transition map}

The construction of \(\theta_{I,J}\) is a Cramer-inverse computation over the
localised chart ring, followed by one application of the universal property of a
localisation. We record the four Mathlib facts it rests on as dependency
anchors.

\begin{lemma}
[Localised element is a unit]
  \label{lem:mathlib_away_algebraMap_isUnit}
  \dcref{custom}
  \lean{IsLocalization.Away.algebraMap_isUnit}
  \mathlibok
  \textit{Provided by Mathlib.}
  Let \(R\) be a commutative ring, \(x \in R\), and \(S\) a localisation of \(R\)
  away from \(x\) (i.e.\ \(S = R[1/x]\)). Then the image of \(x\) under the
  structure map \(R \to S\) is a unit of \(S\).
\end{lemma}

\begin{lemma}
[Left Cramer inverse]
  \label{lem:mathlib_nonsing_inv_mul}
  \dcref{custom}
  \lean{Matrix.nonsing_inv_mul}
  \mathlibok
  \textit{Provided by Mathlib.}
  For a square matrix \(A\) over a commutative ring with \(\det A\) a unit, the
  nonsingular inverse \(A^{-1}\) satisfies \(A^{-1} A = 1\).
\end{lemma}

\begin{lemma}
[Right Cramer inverse]
  \label{lem:mathlib_mul_nonsing_inv}
  \dcref{custom}
  \lean{Matrix.mul_nonsing_inv}
  \mathlibok
  \textit{Provided by Mathlib.}
  For a square matrix \(A\) over a commutative ring with \(\det A\) a unit, the
  nonsingular inverse \(A^{-1}\) satisfies \(A A^{-1} = 1\).
\end{lemma}

\begin{lemma}
[Universal property of an away-localisation]
  \label{lem:mathlib_away_lift}
  \dcref{custom}
  \lean{IsLocalization.Away.lift}
  \mathlibok
  \textit{Provided by Mathlib.}
  Let \(R\) be a commutative ring, \(x \in R\), and \(S = R[1/x]\) a localisation
  of \(R\) away from \(x\). Given any ring homomorphism \(g : R \to P\) that
  sends \(x\) to a unit of \(P\), there is a unique ring homomorphism
  \(S \to P\) extending \(g\) along \(R \to S\).
\end{lemma}

\subsection*{The transition map by matrix algebra}

We now build \(\theta_{I,J}\) in seven steps over the localised chart ring
\(R^I_J := \mathbb{Z}[X^I, 1/P^I_J]\), splitting at the matrix-algebra seams. The
ambient chart ring \(\mathbb{Z}[X^I]\) is the polynomial ring of
\cref{def:gr_affine_chart} in the \(d(r-d)\) free variables \(x^I_{p,q}\)
(\(p \in \{1,\dots,d\}\), \(q \notin I\)).

\begin{definition}[The universal matrix \(X^I\)]\leanok
  \label{def:gr_universal_matrix}
  \dcref{custom}
  \lean{AlgebraicGeometry.Grassmannian.universalMatrix}
  \uses{def:gr_affine_chart}
  % SOURCE: \cite{nitsure-hilbert-quot}, "Construction of Hilbert and Quot Schemes", \S 1,
  % sub-section "Construction by gluing together affine patches" (read from
  % references/nitsure-hilbert-quot-src/nitsure-hilbert-quot.tex, L807-L821).
  % SOURCE QUOTE:
  % "For any subset $I\subset \{1,\ldots, r\}$ with $\#(I) =d$,
  %  consider the $d\times r$ matrix $X^I$ whose $I$ the minor
  %  $X^I_I$ is the $d\times d$ identity matrix $1_{d\times d}$, while the remaining
  %  entries of $X^I$ are independent variables $x^I_{p,q}$ over $\Z$."
  \textit{Source: \cite{nitsure-hilbert-quot}, \S 1.}
  Let \(I \subseteq \{1,\dots,r\}\) with \(\#(I) = d\). The \emph{universal
  matrix} \(X^I\) is the \(d \times r\) matrix over the chart ring
  \(\mathbb{Z}[X^I]\) whose \(I\)-minor \(X^I_I\) is the \(d \times d\) identity
  \(1_{d\times d}\) and whose remaining entries are the free indeterminates
  \(x^I_{p,q}\) (\(1 \le p \le d\), \(q \notin I\)). It is the universal point of
  the chart \(U^I\) (\cref{def:gr_affine_chart}): a ring map
  \(\mathbb{Z}[X^I] \to A\) is the same as a \(d \times r\) matrix over \(A\)
  with identity \(I\)-minor.
\end{definition}

\begin{definition}[The minor determinant \(P^I_J\)]\leanok
  \label{def:gr_minor_det}
  \dcref{custom}
  \lean{AlgebraicGeometry.Grassmannian.minorDet}
  \uses{def:gr_universal_matrix}
  % SOURCE: \cite{nitsure-hilbert-quot}, \S 1, sub-section "Construction by gluing together
  % affine patches" (read from
  % references/nitsure-hilbert-quot-src/nitsure-hilbert-quot.tex, L822-L836).
  % SOURCE QUOTE:
  % "For any $J\subset \{1,\ldots, r\}$ with $\#(J) =d$, let
  %  $P^I_J = \det(X^I_J) \in \Z[X^I]$
  %  where $X^I_J$ is the $J$ th minor of $X^I$."
  \textit{Source: \cite{nitsure-hilbert-quot}, \S 1.}
  For \(J \subseteq \{1,\dots,r\}\) with \(\#(J) = d\), let \(X^I_J\) denote the
  \(J\)-th minor of the universal matrix \(X^I\) (\cref{def:gr_universal_matrix})
  -- the \(d \times d\) submatrix whose columns are those indexed by \(J\). Set
  \[
    P^I_J \;:=\; \det\bigl(X^I_J\bigr) \;\in\; \mathbb{Z}[X^I].
  \]
  Then \(R^I_J := \mathbb{Z}[X^I, 1/P^I_J]\) (the localisation of the chart ring
  away from \(P^I_J\)) is the ring of the principal open subscheme
  \(U^I_J = \Spec \mathbb{Z}[X^I, 1/P^I_J] \subseteq U^I\) on which \(P^I_J\) is
  invertible. One has \(P^I_I = \det(1_{d\times d}) = 1\), so \(R^I_I =
  \mathbb{Z}[X^I]\) and \(U^I_I = U^I\).
\end{definition}

\begin{definition}[The localised \(J\)-minor \(X^I_J\) over \(R^I_J\)]\leanok
  \label{def:gr_universal_minor}
  \dcref{custom}
  \lean{AlgebraicGeometry.Grassmannian.universalMinor}
  \uses{def:gr_universal_matrix, def:gr_minor_det}
  % SOURCE: \cite{nitsure-hilbert-quot}, \S 1, sub-section "Construction by gluing together
  % affine patches" (read from
  % references/nitsure-hilbert-quot-src/nitsure-hilbert-quot.tex, L822-L836).
  % SOURCE QUOTE:
  % "Let $U^I_J = \spec \Z[X^I, 1/P^I_J]$
  %  the open subscheme of $U^I$ where $P^I_J$ is invertible."
  \textit{Source: \cite{nitsure-hilbert-quot}, \S 1.}
  Write \(X^I_J \in \operatorname{Mat}_{d \times d}(R^I_J)\) for the \(J\)-minor of
  the universal matrix \(X^I\) (\cref{def:gr_universal_matrix}) with its entries
  pushed forward along the structure map \(\mathbb{Z}[X^I] \to R^I_J\) into the
  localised chart ring \(R^I_J = \mathbb{Z}[X^I, 1/P^I_J]\)
  (\cref{def:gr_minor_det}). By construction \(\det(X^I_J)\) is the image of
  \(P^I_J\) in \(R^I_J\).
\end{definition}

\begin{lemma}[The localised minor determinant is a unit]\leanok
  \label{lem:gr_minorDet_unit}
  \dcref{custom}
  \lean{AlgebraicGeometry.Grassmannian.isUnit_det_universalMinor}
  \uses{def:gr_universal_minor, def:gr_minor_det, lem:mathlib_away_algebraMap_isUnit}
  % SOURCE: \cite{nitsure-hilbert-quot}, \S 1, sub-section "Construction by gluing together
  % affine patches" (read from
  % references/nitsure-hilbert-quot-src/nitsure-hilbert-quot.tex, L822-L836).
  % SOURCE QUOTE:
  % "the open subscheme of $U^I$ where $P^I_J$ is invertible. This means
  %  the $d\times d$-matrix $X^I_J$ admits an inverse $(X^I_J)^{-1}$
  %  on $U^I_J$."
  \textit{Source: \cite{nitsure-hilbert-quot}, \S 1.}
  The determinant of the localised \(J\)-minor \(X^I_J\)
  (\cref{def:gr_universal_minor}) is a unit of \(R^I_J\):
  \[
    \det\bigl(X^I_J\bigr) \in (R^I_J)^{\times}.
  \]
\end{lemma}
\begin{proof}
  \uses{def:gr_universal_minor, def:gr_minor_det, lem:mathlib_away_algebraMap_isUnit}
  \leanok
  By \cref{def:gr_universal_minor}, \(\det(X^I_J)\) is the image of \(P^I_J\)
  (\cref{def:gr_minor_det}) under the structure map
  \(\mathbb{Z}[X^I] \to R^I_J\). Since \(R^I_J\) is the localisation away from
  \(P^I_J\), that image is a unit by
  \cref{lem:mathlib_away_algebraMap_isUnit}.
\end{proof}

\begin{definition}[The Cramer inverse \((X^I_J)^{-1}\)]\leanok
  \label{def:gr_universalMinorInv}
  \dcref{custom}
  \lean{AlgebraicGeometry.Grassmannian.universalMinorInv}
  \uses{def:gr_universal_minor}
  % SOURCE: \cite{nitsure-hilbert-quot}, \S 1, sub-section "Construction by gluing together
  % affine patches" (read from
  % references/nitsure-hilbert-quot-src/nitsure-hilbert-quot.tex, L822-L836).
  % SOURCE QUOTE:
  % "the $d\times d$-matrix $X^I_J$ admits an inverse $(X^I_J)^{-1}$
  %  on $U^I_J$."
  \textit{Source: \cite{nitsure-hilbert-quot}, \S 1.}
  Let \((X^I_J)^{-1} \in \operatorname{Mat}_{d \times d}(R^I_J)\) be the
  nonsingular (Cramer) inverse of the localised \(J\)-minor \(X^I_J\)
  (\cref{def:gr_universal_minor}). Its entries lie in \(R^I_J\) by Cramer's rule.
\end{definition}

\begin{lemma}[The Cramer inverse is a two-sided inverse]\leanok
  \label{lem:gr_universalMinorInv_identities}
  \dcref{custom}
  \lean{AlgebraicGeometry.Grassmannian.universalMinorInv_mul_cancel}
  \uses{def:gr_universalMinorInv, lem:gr_minorDet_unit, lem:mathlib_nonsing_inv_mul, lem:mathlib_mul_nonsing_inv}
  Because \(\det(X^I_J)\) is a unit of \(R^I_J\) (\cref{lem:gr_minorDet_unit}),
  the Cramer inverse \((X^I_J)^{-1}\) (\cref{def:gr_universalMinorInv}) is a
  genuine two-sided inverse:
  \[
    (X^I_J)^{-1}\, X^I_J = 1_{d\times d}
    \qquad\text{and}\qquad
    X^I_J\, (X^I_J)^{-1} = 1_{d\times d}.
  \]
\end{lemma}
\begin{proof}
  \uses{def:gr_universalMinorInv, lem:gr_minorDet_unit, lem:mathlib_nonsing_inv_mul, lem:mathlib_mul_nonsing_inv}
  \leanok
  Immediate from \cref{lem:gr_minorDet_unit} together with
  \cref{lem:mathlib_nonsing_inv_mul} (left identity) and
  \cref{lem:mathlib_mul_nonsing_inv} (right identity), both of which apply once
  \(\det(X^I_J)\) is known to be a unit.
\end{proof}

\begin{definition}[The image matrix \(M = (X^I_J)^{-1} X^I\)]\leanok
  \label{def:gr_image_matrix}
  \dcref{custom}
  \lean{AlgebraicGeometry.Grassmannian.imageMatrix}
  \uses{def:gr_universalMinorInv, def:gr_universal_matrix}
  % SOURCE: \cite{nitsure-hilbert-quot}, \S 1, sub-section "Construction by gluing together
  % affine patches" (read from
  % references/nitsure-hilbert-quot-src/nitsure-hilbert-quot.tex, L822-L836).
  % SOURCE QUOTE:
  % "The images of the variables
  %  $x^J_{p,q}$ are given by the entries of the matrix formula
  %  $\theta_{I,J}(X^J) = (X^I_J) ^{-1}\,X^I$."
  \textit{Source: \cite{nitsure-hilbert-quot}, \S 1.}
  Define the \(d \times r\) matrix over \(R^I_J\)
  \[
    M \;:=\; (X^I_J)^{-1}\, X^I \;\in\; \operatorname{Mat}_{d \times r}(R^I_J),
  \]
  the product of the Cramer inverse (\cref{def:gr_universalMinorInv}) with the
  universal matrix \(X^I\) (\cref{def:gr_universal_matrix}) base-changed to
  \(R^I_J\). Its entries are the prospective images of the indeterminates
  \(x^J_{p,q}\) under \(\theta_{I,J}\). Its \(J\)-minor is \((X^I_J)^{-1} X^I_J =
  1_{d\times d}\) (\cref{lem:gr_universalMinorInv_identities}), so \(M\) is a
  point of the chart \(U^J\); and its \(I\)-minor is \((X^I_J)^{-1} X^I_I =
  (X^I_J)^{-1}\).
\end{definition}

\begin{definition}[{The pre-localisation hom \(\tilde\theta_{I,J} : \mathbb{Z}[X^J] \to R^I_J\)}]\leanok
  \label{def:gr_transition_pre}
  \dcref{custom}
  \lean{AlgebraicGeometry.Grassmannian.transitionPreMap}
  \uses{def:gr_image_matrix}
  % SOURCE: \cite{nitsure-hilbert-quot}, \S 1, sub-section "Construction by gluing together
  % affine patches" (read from
  % references/nitsure-hilbert-quot-src/nitsure-hilbert-quot.tex, L822-L836).
  % SOURCE QUOTE:
  % "a ring homomorphism
  %  $\theta_{I,J} : \Z[X^J, 1/P^J_I] \to \Z[X^I, 1/P^I_J]$ is
  %  defined as follows. The images of the variables
  %  $x^J_{p,q}$ are given by the entries of the matrix formula
  %  $\theta_{I,J}(X^J) = (X^I_J) ^{-1}\,X^I$."
  \textit{Source: \cite{nitsure-hilbert-quot}, \S 1.}
  Let \(\tilde\theta_{I,J} : \mathbb{Z}[X^J] \to R^I_J\) be the
  \(\mathbb{Z}\)-algebra homomorphism out of the chart ring of \(J\) determined,
  via the universal property of the polynomial ring, by sending each free
  indeterminate \(x^J_{p,q}\) to the \((p,q)\)-entry of the image matrix \(M =
  (X^I_J)^{-1} X^I\) (\cref{def:gr_image_matrix}); equivalently
  \(\tilde\theta_{I,J}(X^J) = (X^I_J)^{-1} X^I\) as a matrix identity (the
  identity \(J\)-block of \(X^J\) matching the identity \(J\)-minor of \(M\)).
\end{definition}

\begin{lemma}[The pre-hom sends \(P^J_I\) to a unit]\leanok
  \label{lem:gr_transition_pre_unit}
  \dcref{custom}
  \lean{AlgebraicGeometry.Grassmannian.isUnit_transitionPreMap_minorDet}
  \uses{def:gr_transition_pre, def:gr_minor_det, lem:gr_universalMinorInv_identities, lem:gr_minorDet_unit}
  % SOURCE: \cite{nitsure-hilbert-quot}, \S 1, sub-section "Construction by gluing together
  % affine patches" (read from
  % references/nitsure-hilbert-quot-src/nitsure-hilbert-quot.tex, L822-L836).
  % SOURCE QUOTE:
  % "In particular, we have $\theta_{I,J}(P^J_I) = 1/P^I_J$,
  %  so the map extends to $\Z[X^J, 1/P^J_I]$."
  \textit{Source: \cite{nitsure-hilbert-quot}, \S 1.}
  The pre-localisation hom \(\tilde\theta_{I,J}\) (\cref{def:gr_transition_pre})
  sends \(P^J_I\) (\cref{def:gr_minor_det}, the \(I\)-minor determinant of
  \(X^J\)) to a unit of \(R^I_J\). Indeed
  \[
    \tilde\theta_{I,J}(P^J_I) = \det\bigl((X^I_J)^{-1} X^I_I\bigr)
      = \det\bigl((X^I_J)^{-1}\bigr) = \det(X^I_J)^{-1} = 1/P^I_J,
  \]
  which is a unit.
\end{lemma}
\begin{proof}
  \uses{def:gr_transition_pre, def:gr_minor_det, lem:gr_universalMinorInv_identities, lem:gr_minorDet_unit}
  \leanok
  By \cref{def:gr_transition_pre}, \(\tilde\theta_{I,J}(X^J) = M = (X^I_J)^{-1}
  X^I\), so \(\tilde\theta_{I,J}\) carries the \(I\)-minor \(X^J_I\) of \(X^J\) to
  the \(I\)-minor of \(M\), which is \((X^I_J)^{-1} X^I_I = (X^I_J)^{-1}\) since
  the \(I\)-minor of \(X^I\) is the identity. Hence \(\tilde\theta_{I,J}(P^J_I) =
  \det((X^I_J)^{-1}) = \det(X^I_J)^{-1}\). As \(\det(X^I_J)\) is a unit
  (\cref{lem:gr_minorDet_unit}), so is its inverse; the matrix-inverse
  determinant identity uses \cref{lem:gr_universalMinorInv_identities}.
\end{proof}

\begin{definition}[Transition map \(\theta_{I,J}\)]\leanok
  \label{def:gr_transition}
  \dcref{custom}
  \lean{AlgebraicGeometry.Grassmannian.transitionMap}
  \uses{def:gr_transition_pre, lem:gr_transition_pre_unit, lem:mathlib_away_lift, def:gr_affine_chart}
  % SOURCE: \cite{nitsure-hilbert-quot}, \S 1, sub-section "Construction by gluing together
  % affine patches" (read from
  % references/nitsure-hilbert-quot-src/nitsure-hilbert-quot.tex, L822-L836).
  % SOURCE QUOTE:
  % "For any $I$ and $J$, a ring homomorphism
  %  $\theta_{I,J} : \Z[X^J, 1/P^J_I] \to \Z[X^I, 1/P^I_J]$ is
  %  defined as follows. ...
  %  In particular, we have $\theta_{I,J}(P^J_I) = 1/P^I_J$,
  %  so the map extends to $\Z[X^J, 1/P^J_I]$."
  \textit{Source: \cite{nitsure-hilbert-quot}, \S 1.}
  Because the pre-localisation hom \(\tilde\theta_{I,J} : \mathbb{Z}[X^J] \to
  R^I_J\) (\cref{def:gr_transition_pre}) sends \(P^J_I\) to a unit of \(R^I_J\)
  (\cref{lem:gr_transition_pre_unit}), the universal property of the
  away-localisation (\cref{lem:mathlib_away_lift}) provides a unique ring
  homomorphism extending it along \(\mathbb{Z}[X^J] \to \mathbb{Z}[X^J, 1/P^J_I]\).
  This extension is the \emph{transition map}
  \[
    \theta_{I,J} : \mathbb{Z}[X^J, 1/P^J_I] \longrightarrow
    \mathbb{Z}[X^I, 1/P^I_J], \qquad \theta_{I,J}(X^J) = (X^I_J)^{-1} X^I .
  \]
  Dually \(\theta_{I,J}\) is the comorphism of a morphism of schemes
  \(U^I_J \to U^J_I\) (\cref{def:gr_affine_chart}), which we also denote
  \(\theta_{I,J}\); it sends the universal matrix \(X^I\) to \((X^I_J)^{-1} X^I\).
\end{definition}

\begin{lemma}[\(\theta_{I,I}\) is the identity]\leanok
  \label{lem:gr_transition_self}
  \dcref{custom}
  \lean{AlgebraicGeometry.Grassmannian.transitionMap_self}
  \uses{def:gr_transition, def:gr_minor_det}
  % SOURCE: \cite{nitsure-hilbert-quot}, \S 1, sub-section "Construction by gluing together
  % affine patches" (read from
  % references/nitsure-hilbert-quot-src/nitsure-hilbert-quot.tex, L838-L848).
  % SOURCE QUOTE:
  % "Note that $\theta_{I,I}$ is identity on $U^I_I = U^I$".
  \textit{Source: \cite{nitsure-hilbert-quot}, \S 1.}
  For every \(I\), \(\theta_{I,I} = \id_{U^I}\): since \(P^I_I = 1\)
  (\cref{def:gr_minor_det}) the minor \(X^I_I = 1_{d\times d}\) is its own
  inverse, so \(\theta_{I,I}(X^I) = (X^I_I)^{-1} X^I = X^I\), i.e.\
  \(\theta_{I,I}\) (\cref{def:gr_transition}) is the identity ring homomorphism
  of \(\mathbb{Z}[X^I, 1/P^I_I] = \mathbb{Z}[X^I]\).
\end{lemma}
\begin{proof}
  \uses{def:gr_transition, def:gr_minor_det}
  \leanok
  When \(J = I\) one has \(P^I_I = \det(1_{d\times d}) = 1\)
  (\cref{def:gr_minor_det}), so \(R^I_I = \mathbb{Z}[X^I]\), the minor \(X^I_I\)
  is the identity, and its Cramer inverse is again the identity. Hence the image
  matrix is \(M = 1 \cdot X^I = X^I\), the pre-hom sends each \(x^I_{p,q}\) to
  itself, and its localisation lift (\cref{def:gr_transition}) is the identity.
\end{proof}

\subsection*{Project-local matrix and minor bookkeeping}

The cocycle computation rests on a handful of elementary matrix-algebra
identities over the (localised) chart rings. They are project-internal helpers
with no external source; each is verified directly in Lean.

\begin{lemma}[Column submatrix of a product]\leanok
  \label{lem:gr_mul_submatrix_col}
  \dcref{custom}
  \lean{AlgebraicGeometry.Grassmannian.mul_submatrix_col}
  Let \(A \in \operatorname{Mat}_{d\times d}(R)\) and
  \(B \in \operatorname{Mat}_{d\times r}(R)\) be matrices over a commutative ring,
  and \(g : \{1,\dots,d\} \to \{1,\dots,r\}\) a column-index map. Then taking the
  column submatrix of the product equals multiplying by the column submatrix of
  the right factor:
  \[
    (A B)_{\bullet\, g} \;=\; A \,\bigl(B_{\bullet\, g}\bigr).
  \]
\end{lemma}
\begin{proof}

\leanok
  Entrywise from the definition of matrix multiplication (reindexing the rows by
  the identity leaves the contraction index untouched). Proved directly in Lean.
\end{proof}

\begin{lemma}[A ring map commutes with the Cramer inverse]\leanok
  \label{lem:gr_map_nonsing_inv}
  \dcref{custom}
  \lean{AlgebraicGeometry.Grassmannian.map_nonsing_inv}
  \uses{lem:mathlib_mul_nonsing_inv}
  Let \(f : R \to S\) be a homomorphism of commutative rings and
  \(A \in \operatorname{Mat}_{n\times n}(R)\) a square matrix with \(\det A\) a
  unit. Writing \(f_\ast\) for the entrywise application of \(f\), the
  nonsingular inverse is natural:
  \[
    \bigl(f_\ast A\bigr)^{-1} \;=\; f_\ast\bigl(A^{-1}\bigr).
  \]
\end{lemma}
\begin{proof}

\leanok
  Since \(f_\ast A \cdot f_\ast(A^{-1}) = f_\ast(A A^{-1}) = f_\ast(1_{n\times n})
  = 1_{n\times n}\) (using \(A A^{-1} = 1\), \cref{lem:mathlib_mul_nonsing_inv}),
  the matrix \(f_\ast(A^{-1})\) is a right inverse of \(f_\ast A\), hence equals
  \((f_\ast A)^{-1}\). Proved directly in Lean.
\end{proof}

\begin{lemma}[Functoriality of entrywise application]\leanok
  \label{lem:gr_map_map_eq_of_comp}
  \dcref{custom}
  \lean{AlgebraicGeometry.Grassmannian.map_map_eq_of_comp}
  Let \(f : R \to A\), \(g : A \to D\) and \(h : R \to D\) be ring
  homomorphisms with \(g \circ f = h\). Then for any matrix \(M\) over \(R\),
  applying \(f\) then \(g\) entrywise equals applying \(h\) entrywise:
  \(g_\ast(f_\ast M) = h_\ast M\).
\end{lemma}
\begin{proof}

\leanok
  Immediate from functoriality of \(\operatorname{Mat}(-)\) (composition of
  entrywise maps). Proved directly in Lean.
\end{proof}

\begin{lemma}[Inverse cancellation in a composite]\leanok
  \label{lem:gr_inv_mul_inv_mul_cancel}
  \dcref{custom}
  \lean{AlgebraicGeometry.Grassmannian.inv_mul_inv_mul_cancel}
  \uses{lem:mathlib_mul_nonsing_inv}
  Let \(A, B \in \operatorname{Mat}_{d\times d}(R)\) and
  \(M \in \operatorname{Mat}_{d\times e}(R)\) over a commutative ring, with
  \(\det A\) a unit. Then
  \[
    \bigl(B^{-1} A\bigr)\,\bigl(A^{-1} M\bigr) \;=\; B^{-1} M.
  \]
\end{lemma}
\begin{proof}

\leanok
  By associativity and \(A A^{-1} = 1_{d\times d}\)
  (\cref{lem:mathlib_mul_nonsing_inv}). Proved directly in Lean.
\end{proof}

\begin{lemma}[The \(I\)-minor of \(X^I\) is the identity]\leanok
  \label{lem:gr_universalMatrix_submatrix_self}
  \dcref{custom}
  \lean{AlgebraicGeometry.Grassmannian.universalMatrix_submatrix_self}
  \uses{def:gr_universal_matrix}
  The \(I\)-minor of the universal matrix \(X^I\)
  (\cref{def:gr_universal_matrix}), with its \(I\)-columns read through the order
  isomorphism \(\{1,\dots,d\} \xrightarrow{\sim} I\), is the \(d\times d\)
  identity: \(X^I_I = 1_{d\times d}\). (This is the normalisation \(P^I_I = 1\)
  defining the chart.)
\end{lemma}
\begin{proof}

\leanok
  By \cref{def:gr_universal_matrix} the \(q\)-column with \(q \in I\) is the
  standard basis vector \(e_p\) where \(p\) is the preimage of \(q\) under the
  order isomorphism \(\{1,\dots,d\} \xrightarrow{\sim} I\); so the \(I\)-minor
  is the identity. Proved directly in Lean.
\end{proof}

\begin{lemma}[The \(J\)-minor of the image matrix is the identity]\leanok
  \label{lem:gr_imageMatrix_submatrix_self}
  \dcref{custom}
  \lean{AlgebraicGeometry.Grassmannian.imageMatrix_submatrix_self}
  \uses{def:gr_image_matrix, lem:gr_mul_submatrix_col, lem:gr_universalMinorInv_identities}
  The \(J\)-minor of the image matrix \(M = (X^I_J)^{-1} X^I\)
  (\cref{def:gr_image_matrix}) is the identity:
  \[
    M_J \;=\; (X^I_J)^{-1}\, X^I_J \;=\; 1_{d\times d}.
  \]
\end{lemma}
\begin{proof}

\leanok
  The \(J\)-column submatrix of \(M = (X^I_J)^{-1} X^I\) is
  \((X^I_J)^{-1}\) times the \(J\)-column submatrix of \(X^I\), i.e.
  \((X^I_J)^{-1} X^I_J\), by \cref{lem:gr_mul_submatrix_col}; this is
  \(1_{d\times d}\) by the left-inverse identity
  (\cref{lem:gr_universalMinorInv_identities}). Proved directly in Lean.
\end{proof}


\begin{lemma}[\(\tilde\theta_{I,J}\) realises the matrix formula]\leanok
  \label{lem:gr_universalMatrix_map_transitionPreMap}
  \dcref{custom}
  \lean{AlgebraicGeometry.Grassmannian.universalMatrix_map_transitionPreMap}
  \uses{def:gr_universal_matrix, def:gr_image_matrix, def:gr_transition_pre, lem:gr_imageMatrix_submatrix_self}
  Applying the pre-localisation hom \(\tilde\theta_{I,J}\)
  (\cref{def:gr_transition_pre}) entrywise to the universal matrix \(X^J\) yields
  the image matrix:
  \[
    (\tilde\theta_{I,J})_\ast\, X^J \;=\; M \;=\; (X^I_J)^{-1} X^I .
  \]
\end{lemma}
\begin{proof}

\leanok
  On a free entry \(x^J_{p,q}\) (\(q \notin J\)) this holds by the definition of
  \(\tilde\theta_{I,J}\) as evaluation at the entries of \(M\)
  (\cref{def:gr_transition_pre}). On a column \(q \in J\) the entry of \(X^J\) is
  an identity-block entry, which \(\tilde\theta_{I,J}\) carries to the
  corresponding entry of \(M_J = 1_{d\times d}\)
  (\cref{lem:gr_imageMatrix_submatrix_self}). Proved directly in Lean.
\end{proof}

\begin{lemma}[Image of a cross minor under \(\tilde\theta_{I,J}\)]\leanok
  \label{lem:gr_transitionPreMap_minorDet}
  \dcref{custom}
  \lean{AlgebraicGeometry.Grassmannian.transitionPreMap_minorDet}
  \uses{def:gr_transition_pre, def:gr_minor_det, def:gr_image_matrix, lem:gr_universalMatrix_map_transitionPreMap}
  For a third size-\(d\) subset \(K\), the pre-hom \(\tilde\theta_{I,J}\) sends the
  minor determinant \(P^J_K\) (\cref{def:gr_minor_det}) to the determinant of the
  \(K\)-minor of the image matrix \(M = (X^I_J)^{-1} X^I\):
  \[
    \tilde\theta_{I,J}\bigl(P^J_K\bigr) \;=\; \det\bigl(M_K\bigr).
  \]
\end{lemma}
\begin{proof}

\leanok
  The hom \(\tilde\theta_{I,J}\) commutes with determinants and with passing to
  a column submatrix, so \(\tilde\theta_{I,J}(P^J_K) =
  \det\bigl((\tilde\theta_{I,J})_\ast X^J\bigr)_K\); apply
  \cref{lem:gr_universalMatrix_map_transitionPreMap}. Proved directly in Lean.
\end{proof}

\begin{lemma}[Naturality of the image matrix]\leanok
  \label{lem:gr_imageMatrix_map_eq}
  \dcref{custom}
  \lean{AlgebraicGeometry.Grassmannian.imageMatrix_map_eq}
  \uses{def:gr_image_matrix, lem:gr_map_map_eq_of_comp, def:gr_universalMinorInv, lem:gr_map_nonsing_inv, lem:gr_minorDet_unit, def:gr_universal_minor}
  Let \(X\) be a size-\(d\) subset and \(\iota : R^I_X \to D\) a ring
  homomorphism lying over the structure map \(R^I \to D\), i.e.\
  \(\iota \circ (R^I \to R^I_X) = (R^I \to D)\). Writing \(Y\) for the universal
  matrix \(X^I\) base-changed to \(D\) and \(Y_X\) for its \(X\)-minor, one has
  \[
    \iota_\ast\bigl((X^I_X)^{-1} X^I\bigr) \;=\; (Y_X)^{-1}\, Y .
  \]
\end{lemma}
\begin{proof}

\leanok
  Distribute \(\iota_\ast\) across the product \(M = (X^I_X)^{-1} X^I\); the
  factor \((X^I_X)^{-1}\) passes through \(\iota_\ast\) as a Cramer inverse
  (\cref{lem:gr_map_nonsing_inv}, using \cref{lem:gr_minorDet_unit}), and both
  factors collapse to base changes over \(D\) by functoriality
  (\cref{lem:gr_map_map_eq_of_comp}) since \(\iota\) lies over \(R^I \to D\).
  Proved directly in Lean.
\end{proof}

\subsection*{Triple-overlap rings and the localised transition maps}

The cocycle condition \(\theta_{I,K} = \theta_{I,J} \circ \theta_{J,K}\) cannot
be stated as a naive composition of the landed transition maps of
\cref{def:gr_transition}: the codomain of \(\theta_{J,K}\) (the ring
\(R^J[1/P^J_K]\)) differs from the domain of \(\theta_{I,J}\) (the ring
\(R^J[1/P^J_I]\)). The identity therefore lives over the \emph{triple-overlap}
rings obtained by inverting both relevant minors in each chart ring:
\[
  S_K := R^K[1/(P^K_I P^K_J)], \quad
  S_J := R^J[1/(P^J_I P^J_K)], \quad
  S_I := R^I[1/(P^I_J P^I_K)].
\]

\begin{definition}[Inclusion of an away-localisation into a double one (left)]\leanok
  \label{def:gr_away_incl_left}
  \dcref{custom}
  \lean{AlgebraicGeometry.Grassmannian.awayInclLeft}
  \uses{lem:mathlib_away_lift, lem:mathlib_away_algebraMap_isUnit}
  For \(x, y \in R\), the canonical inclusion
  \(\iota^{\mathrm{L}}_{x,y} : R[1/x] \to R[1/(xy)]\) inverting the extra factor
  \(y\), defined as the away-localisation lift (\cref{lem:mathlib_away_lift}) of
  the structure map \(R \to R[1/(xy)]\) along \(x\) — well-defined because \(x\)
  is a unit of \(R[1/(xy)]\) (a factor of the unit \(xy\),
  \cref{lem:mathlib_away_algebraMap_isUnit}).
\end{definition}

\begin{definition}[Inclusion of an away-localisation into a double one (right)]\leanok
  \label{def:gr_away_incl_right}
  \dcref{custom}
  \lean{AlgebraicGeometry.Grassmannian.awayInclRight}
  \uses{lem:mathlib_away_lift, lem:mathlib_away_algebraMap_isUnit}
  For \(x, y \in R\), the canonical inclusion
  \(\iota^{\mathrm{R}}_{x,y} : R[1/y] \to R[1/(xy)]\) inverting the extra factor
  \(x\), defined as the away-localisation lift (\cref{lem:mathlib_away_lift}) of
  \(R \to R[1/(xy)]\) along \(y\) — well-defined since \(y\) is a unit of
  \(R[1/(xy)]\) (\cref{lem:mathlib_away_algebraMap_isUnit}).
\end{definition}

\begin{lemma}[Left inclusion lies over \(R\)]\leanok
  \label{lem:gr_awayInclLeft_comp_algebraMap}
  \dcref{custom}
  \lean{AlgebraicGeometry.Grassmannian.awayInclLeft_comp_algebraMap}
  \uses{def:gr_away_incl_left}
  The inclusion \(\iota^{\mathrm{L}}_{x,y}\) (\cref{def:gr_away_incl_left})
  intertwines the two structure maps:
  \(\iota^{\mathrm{L}}_{x,y} \circ (R \to R[1/x]) = (R \to R[1/(xy)])\).
\end{lemma}
\begin{proof}

\leanok
  By the defining property of the away-localisation lift
  (\cref{lem:mathlib_away_lift}). Proved directly in Lean.
\end{proof}

\begin{lemma}[Right inclusion lies over \(R\)]\leanok
  \label{lem:gr_awayInclRight_comp_algebraMap}
  \dcref{custom}
  \lean{AlgebraicGeometry.Grassmannian.awayInclRight_comp_algebraMap}
  \uses{def:gr_away_incl_right}
  The inclusion \(\iota^{\mathrm{R}}_{x,y}\) (\cref{def:gr_away_incl_right})
  intertwines the two structure maps:
  \(\iota^{\mathrm{R}}_{x,y} \circ (R \to R[1/y]) = (R \to R[1/(xy)])\).
\end{lemma}
\begin{proof}

\leanok
  By the defining property of the away-localisation lift
  (\cref{lem:mathlib_away_lift}). Proved directly in Lean.
\end{proof}

\begin{lemma}[Left factor is a unit in the double localisation]\leanok
  \label{lem:gr_isUnit_algebraMap_away_left}
  \dcref{custom}
  \lean{AlgebraicGeometry.Grassmannian.isUnit_algebraMap_away_left}
  \uses{lem:mathlib_away_algebraMap_isUnit}
  For \(x, y \in R\), the image of \(x\) in \(R[1/(xy)]\) is a unit.
\end{lemma}
\begin{proof}

\leanok
  \(xy\) is a unit of \(R[1/(xy)]\)
  (\cref{lem:mathlib_away_algebraMap_isUnit}); a factor of a unit is a unit.
  Proved directly in Lean.
\end{proof}

\begin{lemma}[Right factor is a unit in the double localisation]\leanok
  \label{lem:gr_isUnit_algebraMap_away_right}
  \dcref{custom}
  \lean{AlgebraicGeometry.Grassmannian.isUnit_algebraMap_away_right}
  \uses{lem:mathlib_away_algebraMap_isUnit}
  For \(x, y \in R\), the image of \(y\) in \(R[1/(xy)]\) is a unit.
\end{lemma}
\begin{proof}

\leanok
  \(xy\) is a unit of \(R[1/(xy)]\)
  (\cref{lem:mathlib_away_algebraMap_isUnit}); a factor of a unit is a unit.
  Proved directly in Lean.
\end{proof}

\begin{lemma}[The cross minor becomes a unit after inclusion]\leanok
  \label{lem:gr_isUnit_incl_transitionPreMap_cross}
  \dcref{custom}
  \lean{AlgebraicGeometry.Grassmannian.isUnit_incl_transitionPreMap_cross}
  \uses{lem:gr_transitionPreMap_minorDet, lem:gr_mul_submatrix_col, lem:gr_universalMinorInv_identities, def:gr_image_matrix}
  Let \(A, B, C\) be size-\(d\) subsets and \(\iota : R^A_B \to D\) a ring
  homomorphism over the structure map \(R^A \to D\) under which the minor
  determinant \(P^A_C\) becomes a unit of \(D\). Then
  \(\iota\bigl(\tilde\theta_{A,B}(P^B_C)\bigr)\) is a unit of \(D\). Indeed
  \[
    \tilde\theta_{A,B}\bigl(P^B_C\bigr)
      = \det\bigl((X^A_B)^{-1} X^A_C\bigr)
      = \det\bigl((X^A_B)^{-1}\bigr)\cdot P^A_C ,
  \]
  a product of two units once \(P^A_C\) is inverted in \(D\).
\end{lemma}
\begin{proof}

\leanok
  By \cref{lem:gr_transitionPreMap_minorDet},
  \(\tilde\theta_{A,B}(P^B_C) = \det(M_C)\) with \(M = (X^A_B)^{-1} X^A\); split
  the \(C\)-minor of the product (\cref{lem:gr_mul_submatrix_col}) and the
  determinant as \(\det((X^A_B)^{-1}) \cdot P^A_C\). The first factor is a unit
  (\cref{lem:gr_universalMinorInv_identities}); the second becomes a unit under
  \(\iota\) by hypothesis. Proved directly in Lean.
\end{proof}

\begin{definition}[Localised transition map \(\Theta_{I,J} : S_J \to S_I\)]\leanok
  \label{def:gr_cocycle_theta_ij}
  \dcref{custom}
  \lean{AlgebraicGeometry.Grassmannian.cocycleΘIJ}
  \uses{lem:mathlib_away_lift, def:gr_away_incl_left, def:gr_transition_pre, lem:gr_transition_pre_unit, lem:gr_isUnit_incl_transitionPreMap_cross, lem:gr_awayInclLeft_comp_algebraMap, lem:gr_isUnit_algebraMap_away_right, def:gr_minor_det}
  The triple-overlap transition map
  \(\Theta_{I,J} : S_J = R^J[1/(P^J_I P^J_K)] \to S_I = R^I[1/(P^I_J P^I_K)]\):
  the away-localisation lift (\cref{lem:mathlib_away_lift}) along the doubly
  inverted minor \(P^J_I P^J_K\) of the pre-hom \(\tilde\theta_{I,J}\)
  (\cref{def:gr_transition_pre}) post-composed with the inclusion
  \(\iota^{\mathrm{L}}\) into \(S_I\) (\cref{def:gr_away_incl_left}). It is
  well-defined because both inverted minors map to units of \(S_I\):
  \(\tilde\theta_{I,J}(P^J_I)\) by \cref{lem:gr_transition_pre_unit} (then
  \cref{lem:gr_awayInclLeft_comp_algebraMap}), and \(\tilde\theta_{I,J}(P^J_K)\)
  by \cref{lem:gr_isUnit_incl_transitionPreMap_cross} (with cross factor
  \(P^I_K\) a unit, \cref{lem:gr_isUnit_algebraMap_away_right}).
\end{definition}

\begin{definition}[Localised transition map \(\Theta_{J,K} : S_K \to S_J\)]\leanok
  \label{def:gr_cocycle_theta_jk}
  \dcref{custom}
  \lean{AlgebraicGeometry.Grassmannian.cocycleΘJK}
  \uses{lem:mathlib_away_lift, def:gr_away_incl_right, def:gr_transition_pre, lem:gr_transition_pre_unit, lem:gr_isUnit_incl_transitionPreMap_cross, lem:gr_awayInclRight_comp_algebraMap, lem:gr_isUnit_algebraMap_away_left, def:gr_minor_det}
  The triple-overlap transition map
  \(\Theta_{J,K} : S_K = R^K[1/(P^K_I P^K_J)] \to S_J = R^J[1/(P^J_I P^J_K)]\):
  the away-localisation lift (\cref{lem:mathlib_away_lift}) of
  \(\tilde\theta_{J,K}\) (\cref{def:gr_transition_pre}) post-composed with the
  inclusion \(\iota^{\mathrm{R}}\) into \(S_J\) (\cref{def:gr_away_incl_right}),
  along \(P^K_I P^K_J\). Well-defined by \cref{lem:gr_transition_pre_unit} and
  \cref{lem:gr_isUnit_incl_transitionPreMap_cross} (cross factor \(P^J_I\) a
  unit, \cref{lem:gr_isUnit_algebraMap_away_left}).
\end{definition}

\begin{definition}[Localised transition map \(\Theta_{I,K} : S_K \to S_I\)]\leanok
  \label{def:gr_cocycle_theta_ik}
  \dcref{custom}
  \lean{AlgebraicGeometry.Grassmannian.cocycleΘIK}
  \uses{lem:mathlib_away_lift, def:gr_away_incl_right, def:gr_transition_pre, lem:gr_transition_pre_unit, lem:gr_isUnit_incl_transitionPreMap_cross, lem:gr_awayInclRight_comp_algebraMap, lem:gr_isUnit_algebraMap_away_left, def:gr_minor_det}
  The triple-overlap transition map
  \(\Theta_{I,K} : S_K = R^K[1/(P^K_I P^K_J)] \to S_I = R^I[1/(P^I_J P^I_K)]\):
  the away-localisation lift (\cref{lem:mathlib_away_lift}) of
  \(\tilde\theta_{I,K}\) (\cref{def:gr_transition_pre}) post-composed with the
  inclusion \(\iota^{\mathrm{R}}\) into \(S_I\) (\cref{def:gr_away_incl_right}),
  along \(P^K_I P^K_J\). Well-defined by \cref{lem:gr_transition_pre_unit} and
  \cref{lem:gr_isUnit_incl_transitionPreMap_cross} (cross factor \(P^I_J\) a
  unit, \cref{lem:gr_isUnit_algebraMap_away_left}).
\end{definition}

\begin{lemma}[Image-matrix identity over the triple overlap]\leanok
  \label{lem:gr_cocycle_imageMatrix_eq}
  \dcref{custom}
  \lean{AlgebraicGeometry.Grassmannian.cocycle_imageMatrix_eq}
  \uses{lem:gr_imageMatrix_map_eq, def:gr_cocycle_theta_ij, def:gr_away_incl_right, def:gr_away_incl_left, lem:gr_awayInclRight_comp_algebraMap, lem:gr_awayInclLeft_comp_algebraMap, lem:gr_universalMatrix_map_transitionPreMap, lem:gr_map_map_eq_of_comp, lem:gr_map_nonsing_inv, lem:gr_mul_submatrix_col, lem:gr_inv_mul_inv_mul_cancel, lem:gr_isUnit_algebraMap_away_left, lem:gr_minorDet_unit, def:gr_image_matrix, def:gr_universal_minor}
  Over the triple-overlap ring \(S_I\), the image matrix of \(\theta_{I,K}\)
  equals \(\Theta_{I,J}\) applied entrywise to the image matrix of
  \(\theta_{J,K}\):
  \[
    \iota^{\mathrm{R}}_\ast\bigl((X^I_K)^{-1} X^I\bigr)
      \;=\;
      \bigl(\Theta_{I,J}\circ\iota^{\mathrm{R}}\bigr)_\ast
        \bigl((X^J_K)^{-1} X^J\bigr) .
  \]
  Both sides reduce to \((Y_K)^{-1} Y\), where \(Y\) is the universal matrix
  \(X^I\) base-changed to \(S_I\) and \(Y_K\) its \(K\)-minor.
\end{lemma}
\begin{proof}

\leanok
  The left side is \((Y_K)^{-1} Y\) by naturality of the image matrix
  (\cref{lem:gr_imageMatrix_map_eq}, via \cref{lem:gr_awayInclRight_comp_algebraMap}).
  For the right side, \(\Theta_{I,J}\circ\iota^{\mathrm{R}}\) carries the chart
  ring \(R^J\) over \(\iota^{\mathrm{L}}\circ\tilde\theta_{I,J}\), so the image
  matrix of \(\theta_{J,K}\) maps to \((M^J_K)^{-1} M^J\) with
  \(M^J = (Y_J)^{-1} Y\) (\cref{lem:gr_universalMatrix_map_transitionPreMap},
  \cref{lem:gr_imageMatrix_map_eq}); the minor \(M^J_K = (Y_J)^{-1} Y_K\)
  (\cref{lem:gr_mul_submatrix_col}) inverts via \((AB)^{-1} = B^{-1}A^{-1}\), and
  the cancellation \((Y_K)^{-1} Y_J \cdot (Y_J)^{-1} Y = (Y_K)^{-1} Y\)
  (\cref{lem:gr_inv_mul_inv_mul_cancel}, using \(\det Y_J\) a unit by
  \cref{lem:gr_isUnit_algebraMap_away_left}) gives \((Y_K)^{-1} Y\). Proved
  directly in Lean.
\end{proof}

\begin{lemma}[Cocycle condition for the transition maps]\leanok
  \label{lem:gr_cocycle}
  \dcref{custom}
  \lean{AlgebraicGeometry.Grassmannian.cocycleCondition}
  \uses{def:gr_transition, def:gr_affine_chart, def:gr_image_matrix, lem:gr_universalMinorInv_identities, lem:gr_transition_self, def:gr_cocycle_theta_ij, def:gr_cocycle_theta_jk, def:gr_cocycle_theta_ik, lem:gr_cocycle_imageMatrix_eq}
  % SOURCE: \cite{nitsure-hilbert-quot}, \S 1, sub-section "Construction by gluing together
  % affine patches" (read from
  % references/nitsure-hilbert-quot-src/nitsure-hilbert-quot.tex, L838-L848).
  % SOURCE QUOTE:
  % "Note that $\theta_{I,I}$ is identity on $U^I_I = U^I$, and
  %  we leave it to the reader to verify that
  %  for any three subsets $I$, $J$ and $K$ of $\{ 1,\ldots, r\}$ of
  %  cardinality $d$, the co-cycle condition
  %  $\theta_{I,K}= \theta_{I,J}\theta_{J,K}$
  %  is satisfied."
  \textit{Source: \cite{nitsure-hilbert-quot}, \S 1.}
  For any three subsets \(I, J, K \subseteq \{1, \dots, r\}\) of cardinality
  \(d\), the transition maps of \cref{def:gr_transition} satisfy the cocycle
  condition. As morphisms of schemes, on the triple overlap inside \(U^I\),
  \[
    \theta_{I,K} \;=\; \theta_{J,K} \circ \theta_{I,J},
  \]
  equivalently, as ring homomorphisms,
  \(\theta_{I,K} = \theta_{I,J} \circ \theta_{J,K}\) (the two read the same
  identity through the order-reversal of \(\Spec\)). Together with
  \(\theta_{I,I} = \id\), this is exactly the compatibility datum required to
  glue the charts.
\end{lemma}

% SOURCE QUOTE PROOF: "Note that $\theta_{I,I}$ is identity on $U^I_I = U^I$, and
% we leave it to the reader to verify that
% for any three subsets $I$, $J$ and $K$ of $\{ 1,\ldots, r\}$ of
% cardinality $d$, the co-cycle condition
% $\theta_{I,K}= \theta_{I,J}\theta_{J,K}$
% is satisfied."  (references/nitsure-hilbert-quot-src/nitsure-hilbert-quot.tex,
% L838-L848.)
\begin{proof}
  \uses{def:gr_transition, def:gr_affine_chart, def:gr_image_matrix, lem:gr_universalMinorInv_identities, lem:gr_transition_self, def:gr_cocycle_theta_ij, def:gr_cocycle_theta_jk, def:gr_cocycle_theta_ik, lem:gr_cocycle_imageMatrix_eq}
  \leanok
  We verify the scheme-morphism form
  \(\theta_{I,K} = \theta_{J,K} \circ \theta_{I,J}\) by a direct matrix
  computation, valid over the open locus of \(U^I\) on which all the minors
  involved are invertible (namely \(P^I_J\) and \(P^I_K\), which suffices since
  the maps are determined by their restriction to a dense open and all rings are
  domains; the computation is an identity of matrices with entries in
  \(\mathbb{Z}[X^I, 1/P^I_J, 1/P^I_K]\)).

  Start from the matrix \(X^I\) (the universal point of \(U^I\)). By
  \cref{def:gr_transition},
  \[
    M^J \;:=\; \theta_{I,J}(X^I) \;=\; (X^I_J)^{-1} X^I,
  \]
  a \(d \times r\) matrix whose \(J\)-minor is the identity, hence a point of
  \(U^J\). Its \(K\)-minor is
  \[
    M^J_K \;=\; (X^I_J)^{-1} X^I_K,
    \qquad\text{so}\qquad
    (M^J_K)^{-1} \;=\; (X^I_K)^{-1} X^I_J,
  \]
  using \((AB)^{-1} = B^{-1} A^{-1}\) for the invertible \(d \times d\) matrices
  \(X^I_J\) and \(X^I_K\). Applying \(\theta_{J,K}\) to \(M^J\) gives
  \[
    \theta_{J,K}(M^J) \;=\; (M^J_K)^{-1} M^J
    \;=\; (X^I_K)^{-1} X^I_J \,\cdot\, (X^I_J)^{-1} X^I
    \;=\; (X^I_K)^{-1} X^I,
  \]
  the middle product \(X^I_J (X^I_J)^{-1} = 1_{d\times d}\) cancelling by
  associativity of matrix multiplication. The right-hand side is precisely
  \(\theta_{I,K}(X^I) = (X^I_K)^{-1} X^I\). Thus
  \(\theta_{J,K} \circ \theta_{I,J}\) and \(\theta_{I,K}\) agree on the
  universal matrix \(X^I\), hence agree as morphisms \(U^I \supseteq U^I_J \cap
  U^I_K \to U^K\). The special case \(K = I\) (resp.\ \(J = I\)) together with
  \(\theta_{I,I} = \id\) shows each \(\theta_{I,J}\) is an isomorphism onto its
  image with inverse \(\theta_{J,I}\), since
  \(\theta_{J,I} \circ \theta_{I,J} = \theta_{I,I} = \id\).
\end{proof}

\section{Scheme-level glue-data layer}
\label{sec:gr_glue_layer}

The transition data of \cref{sec:gr_transition} lives at the level of rings. To
feed it to a scheme-gluing machine we must re-express each piece as a morphism in
the category of schemes: the charts \(U^I\), the principal-open overlaps
\(U^I_J\), their inclusions, and the transitions \(\theta_{I,J}\) all become
schemes and scheme morphisms by applying the contravariant functor \(\Spec\).
The lemmas of this section record those translations together with the two
genuinely geometric facts needed to assemble a glue datum: that each inclusion
\(U^I_J \hookrightarrow U^I\) is an open immersion, and that the fibre product of
two principal opens of \(\Spec A\) is again a principal open. They are
project-internal bridge results with no external source beyond the gluing
construction of \cref{def:gr_glued_scheme}.

\subsection*{Mathlib inputs for the scheme-level layer}

\begin{lemma}
[\(\Spec\) of an away-localisation is an open immersion]
  \label{lem:mathlib_specMap_localizationAway_isOpenImmersion}
  \dcref{custom}
  \lean{AlgebraicGeometry.Scheme.isOpenImmersion_SpecMap_localizationAway}
  \mathlibok
  \textit{Provided by Mathlib.}
  Let \(R\) be a commutative ring and \(f \in R\). The morphism
  \(\Spec R[1/f] \to \Spec R\) induced by the structure map \(R \to R[1/f]\) of
  the away-localisation is an open immersion (onto the basic open
  \(D(f) \subseteq \Spec R\)).
\end{lemma}

\begin{lemma}
[Fibre product of two affine spectra is \(\Spec\) of the tensor product]
  \label{lem:mathlib_pullbackSpecIso}
  \dcref{custom}
  \lean{AlgebraicGeometry.pullbackSpecIso}
  \mathlibok
  \textit{Provided by Mathlib.}
  Let \(R\) be a commutative ring and \(S\), \(T\) two \(R\)-algebras. The fibre
  product of \(\Spec S \to \Spec R \leftarrow \Spec T\) is canonically isomorphic
  to the spectrum of the tensor product:
  \[
    \Spec S \times_{\Spec R} \Spec T \;\cong\; \Spec\bigl(S \otimes_R T\bigr).
  \]
\end{lemma}

\begin{lemma}
[An away-localisation of a tensor product is a double away-localisation]
  \label{lem:mathlib_isLocalization_away_mul}
  \dcref{custom}
  \lean{IsLocalization.Away.mul'}
  \mathlibok
  \textit{Provided by Mathlib.}
  Let \(R\) be a commutative ring, \(x, y \in R\), and \(R[1/x]\) the localisation
  away from \(x\). Then \(R[1/x] \otimes_R R[1/y]\) is a localisation of \(R\)
  away from \(x y\); equivalently it realises \(R[1/(xy)]\).
\end{lemma}

\begin{lemma}
[Two away-localisations of the same element are canonically isomorphic]
  \label{lem:mathlib_isLocalization_algEquiv}
  \dcref{custom}
  \lean{IsLocalization.algEquiv}
  \mathlibok
  \textit{Provided by Mathlib.}
  Let \(R\) be a commutative ring and \(M \subseteq R\) a submonoid. Any two
  localisations of \(R\) at \(M\) are uniquely \(R\)-algebra isomorphic; in
  particular, given two rings each presented as \(R[1/x]\) there is a canonical
  \(R\)-algebra isomorphism between them.
\end{lemma}

\subsection*{Charts, overlaps, and inclusions as schemes}

\begin{theorem}[The diagonal minor is the unit]\leanok
  \label{lem:gr_minorDet_self}
  \dcref{custom}
  \lean{AlgebraicGeometry.Grassmannian.minorDet_self}
  \uses{def:gr_minor_det, lem:gr_universalMatrix_submatrix_self}
  For every size-\(d\) subset \(I\), the \(I\)-minor determinant of the universal
  matrix \(X^I\) is the unit of the chart ring:
  \[
    P^I_I \;=\; 1 \;\in\; R^I.
  \]
\end{theorem}
\begin{proof}
  \uses{def:gr_minor_det, lem:gr_universalMatrix_submatrix_self}
  \leanok
  The \(I\)-minor \(X^I_I\) is the \(d \times d\) identity
  (\cref{lem:gr_universalMatrix_submatrix_self}), and \(\det(1_{d\times d}) = 1\).
  Proved directly in Lean.
\end{proof}

\begin{definition}[The overlap scheme \(U^I_J\)]\leanok
  \label{def:gr_chart_overlap}
  \dcref{custom}
  \lean{AlgebraicGeometry.Grassmannian.chartOverlap}
  \uses{def:gr_minor_det, def:gr_affine_chart}
  The principal-open overlap, as a scheme, is the affine spectrum of the
  away-localisation of the chart ring \(R^I\) (\cref{def:gr_affine_chart}) at the
  minor determinant \(P^I_J\) (\cref{def:gr_minor_det}):
  \[
    U^I_J \;:=\; \Spec R^I[1/P^I_J].
  \]
  It is the \(V\)-object of the Grassmannian glue datum.
\end{definition}

\begin{definition}[The overlap inclusion \(U^I_J \to U^I\)]\leanok
  \label{def:gr_chart_incl}
  \dcref{custom}
  \lean{AlgebraicGeometry.Grassmannian.chartIncl}
  \uses{def:gr_chart_overlap, def:gr_affine_chart}
  The canonical inclusion of the overlap into the chart is the \(\Spec\) of the
  structure map \(R^I \to R^I[1/P^I_J]\):
  \[
    U^I_J \;=\; \Spec R^I[1/P^I_J] \longrightarrow \Spec R^I \;=\; U^I .
  \]
  It is the \(f\)-field of the Grassmannian glue datum.
\end{definition}

\begin{lemma}[The overlap inclusion is an open immersion]\leanok
  \label{lem:gr_chartIncl_isOpenImmersion}
  \dcref{custom}
  \lean{AlgebraicGeometry.Grassmannian.chartIncl_isOpenImmersion}
  \uses{def:gr_chart_incl, lem:mathlib_specMap_localizationAway_isOpenImmersion}
  The inclusion \(U^I_J \hookrightarrow U^I\) (\cref{def:gr_chart_incl}) is an
  open immersion, identifying \(U^I_J\) with the basic open \(D(P^I_J)\) of
  \(U^I\).
\end{lemma}
\begin{proof}
  \uses{def:gr_chart_incl, lem:mathlib_specMap_localizationAway_isOpenImmersion}
  \leanok
  This is exactly the away-localisation open-immersion fact
  (\cref{lem:mathlib_specMap_localizationAway_isOpenImmersion}) applied to
  \(f = P^I_J\). Proved directly in Lean.
\end{proof}

\begin{lemma}[The self-overlap inclusion is an isomorphism]\leanok
  \label{lem:gr_chartIncl_self_isIso}
  \dcref{custom}
  \lean{AlgebraicGeometry.Grassmannian.chartIncl_self_isIso}
  \uses{def:gr_chart_incl, lem:gr_minorDet_self}
  For every \(I\), the self-inclusion \(U^I_I \to U^I\) (\cref{def:gr_chart_incl})
  is an isomorphism. This is the \(f_{\mathrm{id}}\)-field of the glue datum.
\end{lemma}
\begin{proof}
  \uses{def:gr_chart_incl, lem:gr_minorDet_self}
  \leanok
  Since \(P^I_I = 1\) is a unit (\cref{lem:gr_minorDet_self}), the structure map
  \(R^I \to R^I[1/P^I_I]\) is the localisation at a unit, hence an isomorphism of
  rings; applying \(\Spec\) gives an isomorphism of schemes. Proved directly in
  Lean.
\end{proof}

\begin{definition}[The scheme-level transition \(U^I_J \to U^J_I\)]\leanok
  \label{def:gr_chart_transition}
  \dcref{custom}
  \lean{AlgebraicGeometry.Grassmannian.chartTransition}
  \uses{def:gr_transition, def:gr_chart_overlap}
  The scheme-level transition is the \(\Spec\) of the ring transition map
  \(\theta_{I,J} : R^J[1/P^J_I] \to R^I[1/P^I_J]\) (\cref{def:gr_transition}):
  \[
    U^I_J \;=\; \Spec R^I[1/P^I_J] \longrightarrow \Spec R^J[1/P^J_I] \;=\; U^J_I .
  \]
  It is the \(t\)-field of the Grassmannian glue datum.
\end{definition}

\begin{lemma}[The self-transition is the identity]\leanok
  \label{lem:gr_chartTransition_self}
  \dcref{custom}
  \lean{AlgebraicGeometry.Grassmannian.chartTransition_self}
  \uses{def:gr_chart_transition, lem:gr_transition_self}
  For every \(I\), the scheme-level self-transition
  \(U^I_I \to U^I_I\) (\cref{def:gr_chart_transition}) is the identity. This is
  the \(t_{\mathrm{id}}\)-field of the glue datum.
\end{lemma}
\begin{proof}
  \uses{def:gr_chart_transition, lem:gr_transition_self}
  \leanok
  By \cref{lem:gr_transition_self} the ring map \(\theta_{I,I}\) is the identity,
  and \(\Spec\) sends the identity ring map to the identity morphism. Proved
  directly in Lean.
\end{proof}

\subsection*{The away-localisation pullback isomorphism}

\begin{definition}[Pullback of two principal opens of \(\Spec A\)]\leanok
  \label{def:gr_away_pullback_iso}
  \dcref{custom}
  \lean{AlgebraicGeometry.Grassmannian.awayPullbackIso}
  \uses{lem:mathlib_pullbackSpecIso, lem:mathlib_isLocalization_away_mul, lem:mathlib_isLocalization_algEquiv}
  Let \(A\) be a commutative ring and \(x, y \in A\). The fibre product of the two
  basic-open inclusions \(\Spec A[1/x] \to \Spec A \leftarrow \Spec A[1/y]\) is
  canonically isomorphic to the spectrum of the doubly-inverted localisation:
  \[
    \Spec A[1/x] \times_{\Spec A} \Spec A[1/y]
      \;\cong\; \Spec A[1/(xy)].
  \]
  The isomorphism is the composite of the tensor-product identification of the
  pullback (\cref{lem:mathlib_pullbackSpecIso}) with the \(\Spec\) of the
  localisation isomorphism
  \(A[1/x] \otimes_A A[1/y] \cong A[1/(xy)]\)
  (\cref{lem:mathlib_isLocalization_away_mul},
  \cref{lem:mathlib_isLocalization_algEquiv}). It is stated over a general base
  ring \(A\) so that it may be instantiated at the triple-overlap rings of the
  Grassmannian glue datum.
\end{definition}

\begin{lemma}[Left leg of the pullback isomorphism]\leanok
  \label{lem:gr_awayPullbackIso_inv_fst}
  \dcref{custom}
  \lean{AlgebraicGeometry.Grassmannian.awayPullbackIso_inv_fst}
  \uses{def:gr_away_pullback_iso, def:gr_away_incl_left}
  Under the identification of \cref{def:gr_away_pullback_iso}, the first
  projection of the pullback corresponds to the left away-inclusion: the
  composite
  \(\Spec A[1/(xy)] \to \Spec A[1/x] \times_{\Spec A} \Spec A[1/y]
    \xrightarrow{\mathrm{pr}_1} \Spec A[1/x]\)
  equals \(\Spec\) of the inclusion \(\iota^{\mathrm{L}}_{x,y}\)
  (\cref{def:gr_away_incl_left}).
\end{lemma}
\begin{proof}
  \uses{def:gr_away_pullback_iso, def:gr_away_incl_left}
  \leanok
  The pullback-projection compatibility of the tensor-product identification
  sends \(\mathrm{pr}_1\) to the left tensor inclusion, which under the
  localisation isomorphism is the structure map \(A[1/x] \to A[1/(xy)]\), i.e.\
  \(\iota^{\mathrm{L}}_{x,y}\); the verification is the universal property of the
  away-localisation. Proved directly in Lean.
\end{proof}

\begin{lemma}[Right leg of the pullback isomorphism]\leanok
  \label{lem:gr_awayPullbackIso_inv_snd}
  \dcref{custom}
  \lean{AlgebraicGeometry.Grassmannian.awayPullbackIso_inv_snd}
  \uses{def:gr_away_pullback_iso, def:gr_away_incl_right}
  Symmetrically, the second projection of the pullback corresponds to the right
  away-inclusion: the composite
  \(\Spec A[1/(xy)] \to \Spec A[1/x] \times_{\Spec A} \Spec A[1/y]
    \xrightarrow{\mathrm{pr}_2} \Spec A[1/y]\)
  equals \(\Spec\) of \(\iota^{\mathrm{R}}_{x,y}\)
  (\cref{def:gr_away_incl_right}).
\end{lemma}
\begin{proof}
  \uses{def:gr_away_pullback_iso, def:gr_away_incl_right}
  \leanok
  As in \cref{lem:gr_awayPullbackIso_inv_fst}, with the right tensor inclusion in
  place of the left. Proved directly in Lean.
\end{proof}

\begin{definition}[{Order-swap isomorphism \(A[1/(xy)] \cong A[1/(yx)]\)}]\leanok
  \label{def:gr_away_mul_comm_equiv}
  \dcref{custom}
  \lean{AlgebraicGeometry.Grassmannian.awayMulCommEquiv}
  \uses{lem:mathlib_isLocalization_algEquiv}
  For \(x, y \in A\), the canonical ring isomorphism
  \[
    A[1/(xy)] \;\cong\; A[1/(yx)]
  \]
  identifying the two away-localisations, which agree because \(xy\) and \(yx\)
  generate the same submonoid of \(A\)
  (\cref{lem:mathlib_isLocalization_algEquiv}). It is inserted into the
  triple-cocycle field of the glue datum to reconcile the opposite product
  orders produced by the two pullback isomorphisms.
\end{definition}

\begin{lemma}[The order-swap lies over the base ring]\leanok
  \label{lem:gr_awayMulCommEquiv_comp_algebraMap}
  \dcref{custom}
  \lean{AlgebraicGeometry.Grassmannian.awayMulCommEquiv_comp_algebraMap}
  \uses{def:gr_away_mul_comm_equiv, lem:mathlib_isLocalization_algEquiv}
  The order-swap isomorphism (\cref{def:gr_away_mul_comm_equiv}) intertwines the
  two structure maps: for \(x, y \in A\),
  \[
    \mathrm{swap}_{x,y} \circ \bigl(A \to A[1/(xy)]\bigr)
      \;=\; \bigl(A \to A[1/(yx)]\bigr).
  \]
\end{lemma}
\begin{proof}
  \uses{def:gr_away_mul_comm_equiv, lem:mathlib_isLocalization_algEquiv}
  \leanok
  Both \(A[1/(xy)]\) and \(A[1/(yx)]\) are localisations of \(A\) at the same
  submonoid (\(xy\) and \(yx\) generate it), and the canonical comparison map
  between two such localisations is an \(A\)-algebra isomorphism
  (\cref{lem:mathlib_isLocalization_algEquiv}); an \(A\)-algebra map commutes with
  the structure maps. Proved directly in Lean.
\end{proof}


\subsection*{The triple-overlap glue field \(t'\)}

The scheme-gluing machine requires, beyond the pairwise transitions \(t_{I,J}\), a
\emph{triple-overlap} field \(t'\) on the relevant fibre products, together with
its compatibility \(t_{\mathrm{fac}}\) with the chart projections. We record both
as the next layer of bridge results; their cocycle coherence is the residual
obligation discharged in \cref{def:gr_glued_scheme}.

\begin{definition}[The triple-overlap transition \(t'\)]\leanok
  \label{def:gr_chart_transition'}
  \dcref{custom}
  \lean{AlgebraicGeometry.Grassmannian.chartTransition'}
  \uses{def:gr_away_pullback_iso, def:gr_cocycle_theta_ij, def:gr_away_mul_comm_equiv}
  The \(t'\)-field of the Grassmannian glue datum is the morphism of schemes
  \[
    t'_{I,J,K} : U^I_J \times_{U^I} U^I_K \longrightarrow U^J_K \times_{U^J} U^J_I ,
  \]
  defined as the composite
  \[
    U^I_J \times_{U^I} U^I_K
      \xrightarrow{\ \sim\ } \Spec S_I
      \xrightarrow{\ \Spec \Theta_{I,J}\ } \Spec S_J
      \xrightarrow{\ \Spec \mathrm{swap}\ } \Spec R^J[1/(P^J_K P^J_I)]
      \xrightarrow{\ \sim\ } U^J_K \times_{U^J} U^J_I ,
  \]
  where the first and last arrows are the pullback identifications of
  \cref{def:gr_away_pullback_iso} -- the source one for the pair \((P^I_J, P^I_K)\)
  over \(R^I\), and the target one (taken inverse) for the pair \((P^J_K, P^J_I)\)
  over \(R^J\) -- the middle arrow is the \(\Spec\) of the localised
  triple-overlap transition \(\Theta_{I,J} : S_J \to S_I\)
  (\cref{def:gr_cocycle_theta_ij}), and \(\mathrm{swap}\) is the \(\Spec\) of the
  order-swap isomorphism \(R^J[1/(P^J_K P^J_I)] \cong S_J = R^J[1/(P^J_I P^J_K)]\)
  (\cref{def:gr_away_mul_comm_equiv}). The order-swap is inserted precisely to
  reconcile the product order \(P^J_K P^J_I\) carried by the target pullback
  identification with the order \(P^J_I P^J_K\) of the codomain \(S_J\) of
  \(\Theta_{I,J}\).
\end{definition}

\begin{lemma}[The \(t'\) ring identity]\leanok
  \label{lem:gr_chartTransition'_ringIdentity}
  \dcref{custom}
  \lean{AlgebraicGeometry.Grassmannian.chartTransition'_ringIdentity}
  \uses{def:gr_cocycle_theta_ij, def:gr_away_mul_comm_equiv, def:gr_away_incl_right, def:gr_away_incl_left, def:gr_transition, lem:gr_awayInclRight_comp_algebraMap, lem:gr_awayMulCommEquiv_comp_algebraMap, lem:mathlib_away_lift}
  Over the triple-overlap rings, the localised transition \(\Theta_{I,J}\)
  (\cref{def:gr_cocycle_theta_ij}) pre-composed with the order-swap and the right
  inclusion equals the left inclusion post-composed with the plain transition
  \(\theta_{I,J}\), as ring homomorphisms \(R^J[1/P^J_I] \to S_I\):
  \[
    \Theta_{I,J} \circ \mathrm{swap} \circ \iota^{\mathrm{R}}_{P^J_K, P^J_I}
      \;=\; \iota^{\mathrm{L}}_{P^I_J, P^I_K} \circ \theta_{I,J} .
  \]
\end{lemma}
\begin{proof}
  \uses{def:gr_cocycle_theta_ij, def:gr_away_mul_comm_equiv, def:gr_away_incl_right, def:gr_away_incl_left, def:gr_transition, lem:gr_awayInclRight_comp_algebraMap, lem:gr_awayMulCommEquiv_comp_algebraMap, lem:mathlib_away_lift}
  \leanok
  Both sides are ring maps out of the away-localisation \(R^J[1/P^J_I]\), so by the
  universal property of the localisation (\cref{lem:mathlib_away_lift}) it suffices
  to check equality after pre-composing with the structure map
  \(R^J \to R^J[1/P^J_I]\), i.e.\ on the chart ring \(R^J\). There the right
  inclusion and the order-swap collapse to structure maps
  (\cref{lem:gr_awayInclRight_comp_algebraMap},
  \cref{lem:gr_awayMulCommEquiv_comp_algebraMap}), and the lift defining
  \(\Theta_{I,J}\) (\cref{def:gr_cocycle_theta_ij}) unwinds to
  \(\iota^{\mathrm{L}} \circ \tilde\theta_{I,J}\)
  (\cref{def:gr_away_incl_left}); the right side likewise unwinds through the lift
  defining \(\theta_{I,J}\) (\cref{def:gr_transition}) to the same
  \(\iota^{\mathrm{L}} \circ \tilde\theta_{I,J}\). Proved directly in Lean.
\end{proof}

\begin{lemma}[The \(t'\) compatibility field \(t_{\mathrm{fac}}\)]\leanok
  \label{lem:gr_chartTransition'_fac}
  \dcref{custom}
  \lean{AlgebraicGeometry.Grassmannian.chartTransition'_fac}
  \uses{def:gr_chart_transition', lem:gr_chartTransition'_ringIdentity, def:gr_away_pullback_iso, lem:gr_awayPullbackIso_inv_fst, lem:gr_awayPullbackIso_inv_snd, def:gr_chart_transition}
  The triple-overlap transition \(t'\) (\cref{def:gr_chart_transition'}) is
  compatible with the projections to the charts:
  \[
    \mathrm{pr}_2 \circ t'_{I,J,K} \;=\; t_{I,J} \circ \mathrm{pr}_1 ,
  \]
  where \(\mathrm{pr}_1\) is the first projection of the source pullback
  \(U^I_J \times_{U^I} U^I_K\), \(\mathrm{pr}_2\) the second projection of the
  target pullback \(U^J_K \times_{U^J} U^J_I\), and \(t_{I,J}\) the scheme-level
  transition (\cref{def:gr_chart_transition}).
\end{lemma}
\begin{proof}
  \uses{def:gr_chart_transition', lem:gr_chartTransition'_ringIdentity, def:gr_away_pullback_iso, lem:gr_awayPullbackIso_inv_fst, lem:gr_awayPullbackIso_inv_snd, def:gr_chart_transition}
  \leanok
  Both fibre products are affine, so the identity may be checked after applying
  \(\Spec\), as an identity of the underlying ring homomorphisms. Rewriting the
  source projection \(\mathrm{pr}_1\) through the pullback identification
  (\cref{lem:gr_awayPullbackIso_inv_fst}) and the target projection
  \(\mathrm{pr}_2\) through \cref{lem:gr_awayPullbackIso_inv_snd} turns the claim
  into the ring identity \cref{lem:gr_chartTransition'_ringIdentity}: the three
  localised pieces of \(t'\) (\cref{def:gr_chart_transition'}) collapse to the
  left inclusion followed by the plain transition \(t_{I,J}\)
  (\cref{def:gr_chart_transition}). Proved directly in Lean.
\end{proof}

\begin{lemma}[The triple-overlap cocycle coherence]\leanok
  \label{lem:gr_chartTransition'_cocycle}
  \dcref{custom}
  \lean{AlgebraicGeometry.Grassmannian.chartTransition'_cocycle}
  \uses{def:gr_chart_transition', lem:gr_cocycle_phi_id, def:gr_away_pullback_iso, def:gr_cocycle_theta_ij, def:gr_away_mul_comm_equiv}
  The triple-overlap transitions \(t'\) (\cref{def:gr_chart_transition'}) satisfy
  the cocycle coherence on the source triple overlap \(U^I_J \times_{U^I} U^I_K\):
  \[
    t'_{I,J,K} \;\circ\; t'_{J,K,I} \;\circ\; t'_{K,I,J} \;=\; \mathrm{id}.
  \]
  This is the \(\mathrm{cocycle}\) field of the Grassmannian glue datum -- the
  categorical realisation of the rotated ring identity \(\Phi = \id\)
  (\cref{lem:gr_cocycle_phi_id}).
\end{lemma}
\begin{proof}
  \uses{def:gr_chart_transition', lem:gr_cocycle_phi_id, def:gr_away_pullback_iso, def:gr_cocycle_theta_ij, def:gr_away_mul_comm_equiv}
  \leanok
  Expand each factor \(t'\) into its constituents
  (\cref{def:gr_chart_transition'}): a source pullback identification, the
  \(\Spec\) of a localised triple-overlap transition \(\Theta\)
  (\cref{def:gr_cocycle_theta_ij}) pre-composed with the \(\Spec\) of an
  order-swap (\cref{def:gr_away_mul_comm_equiv}), and an inverse target pullback
  identification. In the threefold composite the two internal pairs
  \(\mathrm{ap}^{-1} \circ \mathrm{ap}\) of pullback identifications
  (\cref{def:gr_away_pullback_iso}) cancel, adjacent factors sharing the same
  identification. There remain the six \(\Spec\)-arrows of the three \(\Theta\)'s
  and three order-swaps; these collapse into a single \(\Spec\) of their composite
  ring endomorphism \(\Phi\). By \cref{lem:gr_cocycle_phi_id} one has
  \(\Phi = \id\), so this middle \(\Spec\) is the identity, and the outer
  conjugating pullback isomorphism cancels against its inverse, leaving the
  identity of the source triple overlap. Proved directly in Lean.
\end{proof}


\begin{lemma}[The rotated triple-overlap ring cocycle identity \(\Phi = \id\)]\leanok
  \label{lem:gr_cocycle_phi_id}
  \dcref{custom}
  \lean{AlgebraicGeometry.Grassmannian.cocyclePhiId}
  \uses{def:gr_cocycle_theta_ij, def:gr_away_mul_comm_equiv, lem:gr_cocycle_imageMatrix_eq,
    lem:gr_cocycle}
  \textit{The ring identity underlying the \(\mathrm{cocycle}\) field of the glue datum.}
  Fix \(d \le r\) and three subsets \(I,J,K \subseteq \{1,\dots,r\}\) of cardinality \(d\). On the
  triple-overlap chart ring \(R_{IJK} := \mathrm{Localization.Away}\,(P^I_J\cdot P^I_K)\)
  the composite ring endomorphism
  \[
    \Phi := \Theta_{I,J,K} \circ \mathrm{swap}_J \circ \Theta_{J,K,I} \circ \mathrm{swap}_K
            \circ \Theta_{K,I,J} \circ \mathrm{swap}_I
  \]
  equals the identity \(\mathrm{RingHom.id}\,R_{IJK}\), where each \(\Theta_{X,Y,Z} =
  \mathtt{cocycleΘIJ}\,d\,r\,X\,Y\,Z\) (\cref{def:gr_cocycle_theta_ij}, taken for the rotated index
  triple) and each \(\mathrm{swap}_X\) is the matching order-swap algebra equivalence
  \(\mathtt{awayMulCommEquiv}\) (\cref{def:gr_away_mul_comm_equiv}).
\end{lemma}
\begin{proof}
  \uses{lem:gr_cocycle_imageMatrix_eq, lem:gr_cocycle}
  \leanok
  Both endomorphisms are ring homomorphisms out of the localization
  \(R_{IJK} = \mathrm{Localization.Away}\,(P^I_J\cdot P^I_K)\), so by Mathlib's
  \(\mathrm{IsLocalization.ringHom\_ext}\) for the submonoid generated by
  \(P^I_J\cdot P^I_K\) it suffices to
  check equality on the chart-ring generators (the images of the universal matrix entries). On those
  generators \(\Phi\) unwinds, via the order-swaps and the \(\Theta\) definitions, to the rotated
  analogue of the matrix cocycle collapse \((Y_K)^{-1}Y\) already proved in
  \cref{lem:gr_cocycle_imageMatrix_eq} — the same image-matrix identity that underlies the scheme-level
  cocycle condition \cref{lem:gr_cocycle}. Reusing that computation for the rotated triple shows each
  generator is fixed, whence \(\Phi = \id\).
\end{proof}

\section{The glued Grassmannian scheme}
\label{sec:gr_glued}

\begin{definition}[The Grassmannian glue datum]\leanok
  \label{def:gr_the_glue_data}
  \dcref{custom}
  \lean{AlgebraicGeometry.Grassmannian.theGlueData}
  \uses{def:gr_affine_chart, def:gr_chart_overlap, def:gr_chart_incl,
    def:gr_chart_transition, def:gr_chart_transition',
    lem:gr_chartIncl_isOpenImmersion, lem:gr_chartIncl_self_isIso,
    lem:gr_chartTransition_self, lem:gr_chartTransition'_fac,
    lem:gr_chartTransition'_cocycle}
  The \emph{Grassmannian glue datum} is the \(\mathtt{Scheme.GlueData}\) bundle
  indexed by the set \(\{\, I \subseteq \{1,\dots,r\} : \#(I) = d \,\}\) of
  size-\(d\) subsets, assembling the data constructed above:
  \begin{itemize}
    \item objects \(U := U^I\), the affine charts (\cref{def:gr_affine_chart});
    \item overlaps \(V := U^I_J\) (\cref{def:gr_chart_overlap});
    \item inclusions \(f := \bigl(U^I_J \hookrightarrow U^I\bigr)\)
      (\cref{def:gr_chart_incl}), which are open immersions
      (\cref{lem:gr_chartIncl_isOpenImmersion}), with the self-inclusion an
      isomorphism (the \(f_{\mathrm{id}}\) field,
      \cref{lem:gr_chartIncl_self_isIso});
    \item transitions \(t := t_{I,J}\) (\cref{def:gr_chart_transition}), with
      \(t_{\mathrm{id}}\) field \cref{lem:gr_chartTransition_self};
    \item the triple-overlap field \(t' := t'_{I,J,K}\)
      (\cref{def:gr_chart_transition'}), with compatibility field
      \(t_{\mathrm{fac}}\) (\cref{lem:gr_chartTransition'_fac}) and \(\mathrm{cocycle}\)
      field (\cref{lem:gr_chartTransition'_cocycle}).
  \end{itemize}
  The remaining structural fields \(f_{\mathrm{mono}}\) (each inclusion is a
  monomorphism) and \(f_{\mathrm{hasPullback}}\) (the relevant fibre products
  exist) are discharged by instance synthesis from the open-immersion property of
  the inclusions. Assembling these named fields into the
  \(\mathtt{Scheme.GlueData}\) record yields the datum from which the glued scheme
  (\cref{def:gr_glued_scheme}) is formed.
\end{definition}

\begin{definition}[The Grassmannian scheme \(\mathrm{Gr}(r,d)\) over \(\mathbb{Z}\)]\leanok
  \label{def:gr_glued_scheme}
  \dcref{custom}
  \lean{AlgebraicGeometry.Grassmannian.scheme}
  \uses{def:gr_the_glue_data, def:gr_affine_chart, def:gr_transition, lem:gr_cocycle,
    def:gr_chart_overlap, def:gr_chart_incl, lem:gr_chartIncl_isOpenImmersion,
    lem:gr_chartIncl_self_isIso, def:gr_chart_transition,
    lem:gr_chartTransition_self, def:gr_away_pullback_iso,
    lem:gr_awayPullbackIso_inv_fst, lem:gr_awayPullbackIso_inv_snd,
    def:gr_away_mul_comm_equiv, def:gr_cocycle_theta_ij,
    def:gr_chart_transition', lem:gr_chartTransition'_fac,
    lem:gr_chartTransition'_ringIdentity, lem:gr_cocycle_imageMatrix_eq,
    lem:gr_cocycle_phi_id, lem:gr_chartTransition'_cocycle}
  % SOURCE: \cite{nitsure-hilbert-quot}, \S 1, sub-section "Construction by gluing together
  % affine patches" (read from
  % references/nitsure-hilbert-quot-src/nitsure-hilbert-quot.tex, L843-L850).
  % SOURCE QUOTE:
  % "Therefore the schemes
  %  $U^I$, as $I$ varies over all the
  %  ${r\choose d}$ different
  %  subsets of  $\{ 1,\ldots, r\}$ of cardinality $d$,
  %  can be glued together by the co-cycle $(\theta_{I,J})$ to form a
  %  finite-type scheme $\grass(r,d)$ over $\Z$. As each $U^I$ is
  %  isomorphic to $\AA^{d(r-d)}_{\Z}$, it follows that
  %  $\grass(r,d) \to \spec \Z$ is smooth of relative dimension $d(r-d)$."
  \textit{Source: \cite{nitsure-hilbert-quot}, \S 1.}
  Gluing the affine charts \(\{U^I\}\) (\cref{def:gr_affine_chart}), one for each
  of the \(\binom{r}{d}\) subsets \(I \subseteq \{1,\dots,r\}\) of cardinality
  \(d\), along the open subschemes \(U^I_J \subseteq U^I\) via the transition
  isomorphisms \(\theta_{I,J} : U^I_J \xrightarrow{\sim} U^J_I\)
  (\cref{def:gr_transition}) -- which satisfy the cocycle condition
  \(\theta_{I,K} = \theta_{J,K} \circ \theta_{I,J}\) and \(\theta_{I,I} = \id\)
  of \cref{lem:gr_cocycle} -- yields a scheme
  \[
    \mathrm{Gr}(r,d) \longrightarrow \Spec \mathbb{Z}.
  \]
  It is of finite type over \(\mathbb{Z}\) (a finite gluing of finitely
  generated \(\mathbb{Z}\)-algebras), and under the gluing each \(U^I\) is an
  open subscheme of \(\mathrm{Gr}(r,d)\) with \(U^I \cap U^J = U^I_J\).

  \emph{Construction of the glue datum.} The gluing is carried out through the
  Mathlib vehicle for scheme gluing, a glue datum in the sense of a category of
  schemes (the structure \(\mathtt{Scheme.GlueData}\), extending the abstract
  glue-datum interface). Its data are supplied as follows. The index set is the
  set \(J := \{\, I \subseteq \{1,\dots,r\} : \#(I) = d \,\}\) of size-\(d\)
  subsets. The objects are the charts \(U^I\) (\cref{def:gr_affine_chart}); the
  pairwise overlaps are the schemes \(U^I_J\) (\cref{def:gr_chart_overlap}); the
  inclusions \(f_{I,J} : U^I_J \to U^I\) are the open immersions
  \(\cref{def:gr_chart_incl}\) (open immersions by
  \cref{lem:gr_chartIncl_isOpenImmersion}, in particular monomorphisms), with the
  identity field \(f_{\mathrm{id}}\) supplied by
  \cref{lem:gr_chartIncl_self_isIso}; the transition isomorphisms
  \(t_{I,J} : U^I_J \to U^J_I\) are the scheme-level transitions
  \(\cref{def:gr_chart_transition}\), with identity field \(t_{\mathrm{id}}\)
  supplied by \cref{lem:gr_chartTransition_self}.

  The triple-overlap field \(t'\) and its first coherence obligation are supplied
  by the two preceding results. The field itself is the morphism
  \(t'_{I,J,K} : U^I_J \times_{U^I} U^I_K \to U^J_K \times_{U^J} U^J_I\) of
  \cref{def:gr_chart_transition'}, and the compatibility field
  \(t_{\mathrm{fac}}\) -- that \(t'\) is compatible with the projections to the
  charts -- is \cref{lem:gr_chartTransition'_fac}, which reduces (both fibre
  products being affine) to the ring identity
  \cref{lem:gr_chartTransition'_ringIdentity} via the leg lemmas
  \cref{lem:gr_awayPullbackIso_inv_fst} and \cref{lem:gr_awayPullbackIso_inv_snd}.

  \emph{Cocycle field.} The \(t'\)-cocycle coherence
  \(t'_{I,J,K} \circ t'_{J,K,I} \circ t'_{K,I,J} = \mathrm{id}\)
  follows from \cref{lem:gr_chartTransition'_cocycle} via the ring identity
  \cref{lem:gr_cocycle_phi_id}. With this field in hand the glue datum is
  complete: it is the structure \(\mathtt{Scheme.GlueData}\) with
  \(U := \texttt{affineChart}\) (\cref{def:gr_affine_chart}),
  \(V := \texttt{chartOverlap}\) (\cref{def:gr_chart_overlap}),
  \(f := \texttt{chartIncl}\) (\cref{def:gr_chart_incl}),
  \(t := \texttt{chartTransition}\) (\cref{def:gr_chart_transition}),
  \(t_{\mathrm{id}} := \texttt{chartTransition\_self}\)
  (\cref{lem:gr_chartTransition_self}),
  \(t' := \texttt{chartTransition'}\) (\cref{def:gr_chart_transition'}),
  \(t_{\mathrm{fac}} := \texttt{chartTransition'\_fac}\)
  (\cref{lem:gr_chartTransition'_fac}), and the cocycle field above. The glued
  scheme \(\mathrm{Gr}(r,d)\) is then the colimit
  \((\text{glue datum}).\mathrm{glued}\) produced by the gluing machine.

  Moreover \(\mathrm{Gr}(r,d) \to \Spec \mathbb{Z}\) is \emph{smooth of relative
  dimension \(d(r-d)\)}. Smoothness and relative dimension are local on the
  source; the charts \(\{U^I\}\) cover \(\mathrm{Gr}(r,d)\), and each
  \(U^I \cong \mathbb{A}^{d(r-d)}_{\mathbb{Z}}\) (\cref{def:gr_affine_chart}) is
  affine space of dimension \(d(r-d)\) over \(\mathbb{Z}\), which is smooth over
  \(\mathbb{Z}\) of relative dimension \(d(r-d)\) (the polynomial ring
  \(\mathbb{Z}[x_1, \dots, x_{d(r-d)}]\) is a smooth \(\mathbb{Z}\)-algebra). As
  the property may be checked on an open cover of the source, the conclusion
  follows.
\end{definition}

\section{Separatedness}
\label{sec:gr_separated}

The surjectivity of the restricted-diagonal comorphism is the ring-theoretic
heart of separatedness. We isolate it through the following sequence of
helper results before stating the keystone lemma.

\begin{lemma}[Two-sided inverse identity \(\tilde\theta_{I,J}(P^J_I)\,P^I_J = 1\)]\leanok
  \label{lem:gr_transitionPreMap_minorDet_swap_mul}
  \dcref{custom}
  \lean{AlgebraicGeometry.Grassmannian.transitionPreMap_minorDet_swap_mul}
  \uses{lem:gr_universalMatrix_map_transitionPreMap, lem:gr_transition_pre_unit, lem:gr_transitionPreMap_minorDet}
  In \(R^I_J = \mathbb{Z}[X^I, 1/P^I_J]\), the pre-localisation transition hom
  \(\tilde\theta_{I,J}\) (\cref{def:gr_transition_pre}) satisfies the explicit
  two-sided inverse identity
  \[
    \tilde\theta_{I,J}\bigl(P^J_I\bigr) \cdot P^I_J \;=\; 1 .
  \]
  This refines the unit assertion of \cref{lem:gr_transition_pre_unit} to a named
  inverse: \(\tilde\theta_{I,J}(P^J_I)\) is the explicit inverse of \(P^I_J\) in
  \(R^I_J\).
\end{lemma}
\begin{proof}
  \uses{lem:gr_universalMatrix_map_transitionPreMap, lem:gr_transition_pre_unit, lem:gr_transitionPreMap_minorDet}
  \leanok
  By \cref{lem:gr_transitionPreMap_minorDet} (taking the third subset \(K = I\)),
  \(\tilde\theta_{I,J}(P^J_I)\) is the determinant of the \(I\)-minor of the image
  matrix \(M = (X^I_J)^{-1} X^I\), which by
  \cref{lem:gr_universalMatrix_map_transitionPreMap} equals
  \(\det\bigl((X^I_J)^{-1} X^I_I\bigr) = \det\bigl((X^I_J)^{-1}\bigr)\). Hence
  \(\tilde\theta_{I,J}(P^J_I) = \det(X^I_J)^{-1}\), and multiplying by \(P^I_J =
  \det(X^I_J)\) gives \(\det(X^I_J)^{-1} \cdot \det(X^I_J) = 1\) by the
  localisation-inverse cancellation for the unit \(\det(X^I_J)\)
  (\cref{lem:gr_transition_pre_unit}).
\end{proof}

\begin{definition}[Restricted-diagonal comorphism \(\delta_{I,J}\)]\leanok
  \label{def:gr_diagonalRingMap}
  \dcref{custom}
  \lean{AlgebraicGeometry.Grassmannian.diagonalRingMap}
  \uses{def:gr_transition_pre}
  The \emph{restricted-diagonal comorphism} on the patch \(U^I \times_{\mathbb{Z}}
  U^J\) is the \(\mathbb{Z}\)-algebra homomorphism
  \[
    \delta_{I,J} : \mathbb{Z}[X^I] \otimes_{\mathbb{Z}} \mathbb{Z}[X^J]
    \longrightarrow R^I_J,
    \qquad
    X^I \otimes 1 \mapsto X^I, \quad
    1 \otimes X^J \mapsto (X^I_J)^{-1} X^I,
  \]
  whose surjectivity is the content of separatedness; this is the \(\delta_{I,J}\)
  of \cref{lem:gr_separated}.
\end{definition}
\begin{proof}
  \uses{def:gr_transition_pre}
  \leanok
  Construct \(\delta_{I,J}\) as the tensor-product lift of the two commuting
  factors: the structure map \(\mathbb{Z}[X^I] = R^I \to R^I_J\) on the left
  factor, and the pre-localisation transition hom \(\tilde\theta_{I,J}\)
  (\cref{def:gr_transition_pre}) on the right factor. Since both land in the
  commutative ring \(R^I_J\), the universal property of the tensor product of
  \(\mathbb{Z}\)-algebras yields a unique \(\mathbb{Z}\)-algebra homomorphism out
  of \(\mathbb{Z}[X^I] \otimes_{\mathbb{Z}} \mathbb{Z}[X^J]\) restricting to these
  factors; by construction its right component is \(\tilde\theta_{I,J}\).
\end{proof}

\begin{lemma}[Left component of \(\delta_{I,J}\)]\leanok
  \label{lem:gr_diagonalRingMap_left}
  \dcref{custom}
  \lean{AlgebraicGeometry.Grassmannian.diagonalRingMap_left}
  \uses{def:gr_diagonalRingMap}
  For \(a \in \mathbb{Z}[X^I]\), the comorphism \(\delta_{I,J}\)
  (\cref{def:gr_diagonalRingMap}) acts on the left factor by the structure map:
  \[
    \delta_{I,J}(a \otimes 1) \;=\; \mathrm{alg}_{R^I \to R^I_J}(a),
  \]
  the image of \(a\) under \(\mathbb{Z}[X^I] = R^I \to R^I_J\).
\end{lemma}
\begin{proof}
  \uses{def:gr_diagonalRingMap}
  \leanok
  Immediate from the defining tensor-product lift of
  \cref{def:gr_diagonalRingMap}: the value on a left elementary tensor \(a \otimes
  1\) is the left factor evaluated at \(a\), i.e.\ the structure map.
\end{proof}

\begin{lemma}[Right component of \(\delta_{I,J}\)]\leanok
  \label{lem:gr_diagonalRingMap_right}
  \dcref{custom}
  \lean{AlgebraicGeometry.Grassmannian.diagonalRingMap_right}
  \uses{def:gr_diagonalRingMap}
  For \(b \in \mathbb{Z}[X^J]\), the comorphism \(\delta_{I,J}\)
  (\cref{def:gr_diagonalRingMap}) acts on the right factor by the pre-localisation
  transition hom:
  \[
    \delta_{I,J}(1 \otimes b) \;=\; \tilde\theta_{I,J}(b)
  \]
  (\cref{def:gr_transition_pre}).
\end{lemma}
\begin{proof}
  \uses{def:gr_diagonalRingMap}
  \leanok
  Immediate from the defining tensor-product lift of
  \cref{def:gr_diagonalRingMap}: the value on a right elementary tensor \(1 \otimes
  b\) is the right factor \(\tilde\theta_{I,J}\) evaluated at \(b\).
\end{proof}

\begin{lemma}[The restricted-diagonal comorphism is surjective]\leanok
  \label{lem:gr_diagonalRingMap_surjective}
  \dcref{custom}
  \lean{AlgebraicGeometry.Grassmannian.diagonalRingMap_surjective}
  \uses{lem:gr_diagonalRingMap_left, lem:gr_diagonalRingMap_right, lem:gr_transitionPreMap_minorDet_swap_mul}
  The comorphism \(\delta_{I,J} : \mathbb{Z}[X^I] \otimes_{\mathbb{Z}}
  \mathbb{Z}[X^J] \to R^I_J\) (\cref{def:gr_diagonalRingMap}) is surjective.
\end{lemma}
\begin{proof}
  \uses{lem:gr_diagonalRingMap_left, lem:gr_diagonalRingMap_right, lem:gr_transitionPreMap_minorDet_swap_mul}
  \leanok
  Let \(z \in R^I_J\). Since \(R^I_J = \mathbb{Z}[X^I, 1/P^I_J]\) is the
  localisation of \(R^I = \mathbb{Z}[X^I]\) away from \(P^I_J\), surjectivity of
  the localisation map provides \(a \in R^I\) and \(n \in \mathbb{N}\) with
  \[
    z \cdot (P^I_J)^n \;=\; \mathrm{alg}_{R^I \to R^I_J}(a).
  \]
  Consider the witness \(a \otimes (P^J_I)^n \in \mathbb{Z}[X^I]
  \otimes_{\mathbb{Z}} \mathbb{Z}[X^J]\). Using multiplicativity together with
  \cref{lem:gr_diagonalRingMap_left} and \cref{lem:gr_diagonalRingMap_right},
  \[
    \delta_{I,J}\bigl(a \otimes (P^J_I)^n\bigr)
      = \mathrm{alg}_{R^I \to R^I_J}(a) \cdot \tilde\theta_{I,J}(P^J_I)^n
      = z \cdot (P^I_J)^n \cdot \tilde\theta_{I,J}(P^J_I)^n
      = z \cdot \bigl(P^I_J \cdot \tilde\theta_{I,J}(P^J_I)\bigr)^n
      = z,
  \]
  where the last equality uses \(P^I_J \cdot \tilde\theta_{I,J}(P^J_I) = 1\) from
  \cref{lem:gr_transitionPreMap_minorDet_swap_mul}. Hence \(z\) lies in the image
  of \(\delta_{I,J}\), so \(\delta_{I,J}\) is surjective.
\end{proof}

\begin{definition}[Source pullback iso \(e_2\) for the diagonal computation]\leanok
  \label{def:gr_pullbackιIso}
  \dcref{custom}
  \lean{AlgebraicGeometry.Grassmannian.pullbackιIso}
  \uses{def:gr_glued_scheme, def:gr_chart_overlap}
  The \emph{source pullback iso} \(e_2\) identifies the categorical fibre product
  of two chart inclusions of the glued scheme (\cref{def:gr_glued_scheme}) with
  the overlap chart:
  \[
    e_2 : \mathrm{pullback}\,(\iota_i)\,(\iota_j) \;\xrightarrow{\;\sim\;}\;
    U^I \cap U^J \;=\; \texttt{chartOverlap}\,(d,r,i,j)
  \]
  (\cref{def:gr_chart_overlap}), where \(\iota_i, \iota_j\) are the open
  immersions of the charts \(U^I, U^J\) into \(\mathrm{Gr}(r,d)\).
\end{definition}
\begin{proof}
  \uses{def:gr_glued_scheme, def:gr_chart_overlap}
  \leanok
  The glue data of \cref{def:gr_glued_scheme} supplies, for each pair \((i,j)\), a
  pullback cone over \(\iota_i, \iota_j\) with apex the overlap chart \(U^I \cap
  U^J\), and this cone is a limit cone. Comparing it with the categorical
  pullback limit cone via the uniqueness-up-to-isomorphism of limit cone points
  yields the desired isomorphism \(e_2\).
\end{proof}

\begin{definition}[Structure morphism \(\mathrm{Gr}(r,d) \to \Spec \mathbb{Z}\)]\leanok
  \label{def:gr_to_specZ}
  \dcref{custom}
  \lean{AlgebraicGeometry.Grassmannian.toSpecZ}
  \uses{def:gr_glued_scheme}
  The \emph{structure morphism} of the glued Grassmannian scheme
  (\cref{def:gr_glued_scheme}) is the canonical map
  \[
    \mathrm{Gr}(r,d) \longrightarrow \Spec \mathbb{Z} .
  \]
  As \(\mathrm{Gr}(r,d)\) lives in the category of schemes over \(\mathbb{Z}\), in
  which \(\Spec \mathbb{Z}\) is the terminal object, this morphism is the unique
  map to that terminal object. It is the genuine base over which separatedness and
  properness of the Grassmannian are formulated.
\end{definition}

\begin{lemma}[The chart map composes to the affine structure map]\leanok
  \label{lem:gr_ι_toSpecZ}
  \dcref{custom}
  \lean{AlgebraicGeometry.Grassmannian.ι_toSpecZ}
  \uses{def:gr_to_specZ, def:gr_glued_scheme}
  For each index \(i = (I, \cdot)\), the chart immersion \(\iota_i : U^I \to
  \mathrm{Gr}(r,d)\) of the glue datum (\cref{def:gr_glued_scheme}) composed with
  the structure morphism (\cref{def:gr_to_specZ}) is the affine structure map of
  the chart:
  \[
    \iota_i \circ (\mathrm{Gr}(r,d) \to \Spec \mathbb{Z})
      \;=\; \Spec(\mathbb{Z} \to R^I),
  \]
  the \(\Spec\) of the canonical ring map \(\mathbb{Z} \to R^I =
  \mathbb{Z}[X^I]\).
\end{lemma}
\begin{proof}
  \uses{def:gr_to_specZ, def:gr_glued_scheme}
  \leanok
  Both sides are morphisms from \(U^I\) into the terminal object \(\Spec
  \mathbb{Z}\), hence they coincide by uniqueness of the terminal map.
\end{proof}

\begin{lemma}[First projection leg of the source pullback iso]\leanok
  \label{lem:gr_pullbackιIso_inv_fst}
  \dcref{custom}
  \lean{AlgebraicGeometry.Grassmannian.pullbackιIso_inv_fst}
  \uses{def:gr_pullbackιIso, def:gr_chart_incl}
  The inverse of the source pullback iso \(e_2\) (\cref{def:gr_pullbackιIso}),
  followed by the first projection of the fibre product, is the chart inclusion:
  \[
    e_2^{-1} \circ \mathrm{pr}_1 \;=\; \iota^I_J : U^I_J \to U^I ,
  \]
  the \(\mathrm{pr}_1\) leg being the projection of
  \(\mathrm{pullback}\,(\iota_i)\,(\iota_j)\) onto \(U^I\)
  (\cref{def:gr_chart_incl}).
\end{lemma}
\begin{proof}
  \uses{def:gr_pullbackιIso, def:gr_chart_incl}
  \leanok
  By construction \(e_2\) is the comparison isomorphism between the categorical
  pullback limit cone and the glue-datum pullback cone over \(\iota_i, \iota_j\).
  The compatibility of a limit-cone comparison isomorphism with the cone legs
  (uniqueness-up-to-isomorphism of limit cone points, applied to the left leg of
  the cospan) identifies \(e_2^{-1} \circ \mathrm{pr}_1\) with the corresponding
  leg of the glue-datum cone, which is the chart inclusion \(\iota^I_J\).
\end{proof}

\begin{lemma}[Second projection leg of the source pullback iso]\leanok
  \label{lem:gr_pullbackιIso_inv_snd}
  \dcref{custom}
  \lean{AlgebraicGeometry.Grassmannian.pullbackιIso_inv_snd}
  \uses{def:gr_pullbackιIso, def:gr_chart_incl, def:gr_chart_transition}
  The inverse of the source pullback iso \(e_2\) (\cref{def:gr_pullbackιIso}),
  followed by the second projection of the fibre product, is the scheme-level
  transition composed with the opposite chart inclusion:
  \[
    e_2^{-1} \circ \mathrm{pr}_2
      \;=\; t_{I,J} \circ \iota^J_I : U^I_J \to U^J_I \to U^J ,
  \]
  the transition \(t_{I,J}\) of \cref{def:gr_chart_transition} followed by the
  chart inclusion of \cref{def:gr_chart_incl}.
\end{lemma}
\begin{proof}
  \uses{def:gr_pullbackιIso, def:gr_chart_incl, def:gr_chart_transition}
  \leanok
  As in \cref{lem:gr_pullbackιIso_inv_fst}, but using the right leg of the cospan:
  the glue-datum pullback cone's second leg is the composite \(t \circ f\) of the
  glue datum (transition followed by inclusion), so the comparison isomorphism
  identifies \(e_2^{-1} \circ \mathrm{pr}_2\) with \(t_{I,J} \circ \iota^J_I\).
\end{proof}

\begin{lemma}[The transition-then-inclusion composite is affine]\leanok
  \label{lem:gr_chartTransition_comp_chartIncl}
  \dcref{custom}
  \lean{AlgebraicGeometry.Grassmannian.chartTransition_comp_chartIncl}
  \uses{def:gr_chart_transition, def:gr_chart_incl, def:gr_transition_pre}
  The composite of the scheme-level transition
  (\cref{def:gr_chart_transition}) with the opposite chart inclusion
  (\cref{def:gr_chart_incl}) is the \(\Spec\) of the pre-localisation transition
  homomorphism \(\tilde\theta_{I,J} : R^J \to R^I[1/P^I_J]\)
  (\cref{def:gr_transition_pre}):
  \[
    t_{I,J} \circ \iota^J_I \;=\; \Spec \tilde\theta_{I,J} :
      U^I_J \to U^J .
  \]
\end{lemma}
\begin{proof}
  \uses{def:gr_chart_transition, def:gr_chart_incl, def:gr_transition_pre}
  \leanok
  Unfolding the two scheme maps, the composite is \(\Spec\) of the ring map
  obtained by precomposing the localised transition \(\theta_{I,J} :
  R^J[1/P^J_I] \to R^I[1/P^I_J]\) with the localisation \(R^J \to R^J[1/P^J_I]\).
  By the universal property of the away-localisation (the lift commutes with the
  localisation map), this composite ring map is exactly the pre-localisation
  transition \(\tilde\theta_{I,J}\); applying \(\Spec\) and the contravariance of
  composition gives the claim.
\end{proof}

\begin{lemma}[The structure morphism is separated (morphism form)]\leanok
  \label{lem:gr_separated_toSpecZ}
  \dcref{custom}
  \lean{AlgebraicGeometry.Grassmannian.isSeparatedToSpecZ}
  \uses{def:gr_to_specZ, def:gr_glued_scheme, def:gr_pullbackιIso,
    def:gr_diagonalRingMap, lem:gr_diagonalRingMap_surjective,
    lem:gr_ι_toSpecZ, lem:gr_pullbackιIso_inv_fst,
    lem:gr_pullbackιIso_inv_snd, lem:gr_chartTransition_comp_chartIncl}
  The structure morphism \(\mathrm{Gr}(r,d) \to \Spec \mathbb{Z}\)
  (\cref{def:gr_to_specZ}) is separated: its diagonal is a closed immersion.
\end{lemma}
\begin{proof}
  \uses{def:gr_to_specZ, def:gr_glued_scheme, def:gr_pullbackιIso,
    def:gr_diagonalRingMap, lem:gr_diagonalRingMap_surjective,
    lem:gr_ι_toSpecZ, lem:gr_pullbackιIso_inv_fst,
    lem:gr_pullbackιIso_inv_snd, lem:gr_chartTransition_comp_chartIncl}
    \leanok
  Separatedness is local on the target of the diagonal, so it may be checked on
  the cover of \(\mathrm{Gr}(r,d) \times_{\mathbb{Z}} \mathrm{Gr}(r,d)\) by the
  patch products \(U^I \times_{\mathbb{Z}} U^J\) (the left--right product of the
  chart open cover with itself). Fix a patch \((i,j)\). On it the restricted
  diagonal is identified, via the source pullback iso \(e_2\)
  (\cref{def:gr_pullbackιIso}) — whose two legs are computed in
  \cref{lem:gr_pullbackιIso_inv_fst} and \cref{lem:gr_pullbackιIso_inv_snd},
  using \cref{lem:gr_ι_toSpecZ} and \cref{lem:gr_chartTransition_comp_chartIncl}
  to read off the affine comorphisms — and the target affine iso (the pullback of
  two affine schemes is the \(\Spec\) of the tensor product), with the affine
  morphism \(\Spec \delta_{I,J}\), where \(\delta_{I,J}\) is the
  restricted-diagonal comorphism of \cref{def:gr_diagonalRingMap}. Since
  \(\delta_{I,J}\) is surjective (\cref{lem:gr_diagonalRingMap_surjective}), the
  morphism \(\Spec \delta_{I,J}\) is a closed immersion. As every patch of the
  diagonal is a closed immersion, the diagonal is a closed immersion and the
  structure morphism is separated.
\end{proof}

\begin{lemma}[\(\mathrm{Gr}(r,d)\) is separated over \(\mathbb{Z}\)]\leanok
  \label{lem:gr_separated}
  \dcref{custom}
  \lean{AlgebraicGeometry.Grassmannian.isSeparated}
  \uses{def:gr_glued_scheme, def:gr_transition, lem:gr_diagonalRingMap_surjective, lem:gr_separated_toSpecZ}
  % SOURCE: \cite{nitsure-hilbert-quot}, \S 1, sub-section "Separatedness" (read from
  % references/nitsure-hilbert-quot-src/nitsure-hilbert-quot.tex, L856-L860).
  % SOURCE QUOTE:
  % "The intersection of the diagonal of $\grass(r,d)$ with $U^I\times U^J$
  %  can be seen to be the closed subscheme $\Delta_{I,J}\subset U^I\times U^J$
  %  defined by entries of the matrix formula $X^J_I\,X^I - X^J =0$,
  %  and so $\grass(r,d)$ is a separated scheme."
  \textit{Source: \cite{nitsure-hilbert-quot}, \S 1.}
  The structure morphism \(\mathrm{Gr}(r,d) \to \Spec \mathbb{Z}\)
  (\cref{def:gr_glued_scheme}) is separated; equivalently, the diagonal
  \(\Delta : \mathrm{Gr}(r,d) \to \mathrm{Gr}(r,d) \times_{\mathbb{Z}}
  \mathrm{Gr}(r,d)\) is a closed immersion.
\end{lemma}

% SOURCE QUOTE PROOF: "The intersection of the diagonal of $\grass(r,d)$ with
% $U^I\times U^J$ can be seen to be the closed subscheme
% $\Delta_{I,J}\subset U^I\times U^J$ defined by entries of the matrix formula
% $X^J_I\,X^I - X^J =0$, and so $\grass(r,d)$ is a separated scheme."
% (references/nitsure-hilbert-quot-src/nitsure-hilbert-quot.tex, L856-L860.)
\begin{proof}
  \uses{def:gr_glued_scheme, def:gr_transition, lem:gr_diagonalRingMap_surjective, lem:gr_separated_toSpecZ}
  \leanok
  The products \(U^I \times_{\mathbb{Z}} U^J\), as \(I, J\) range over the
  size-\(d\) subsets, form an affine open cover of \(\mathrm{Gr}(r,d)
  \times_{\mathbb{Z}} \mathrm{Gr}(r,d)\). Separatedness may be checked on this
  cover: it suffices to show that for each pair \((I,J)\) the preimage
  \(\Delta^{-1}(U^I \times_{\mathbb{Z}} U^J)\) maps by a closed immersion into
  \(U^I \times_{\mathbb{Z}} U^J\), i.e.\ that the corresponding ring
  homomorphism is surjective.

  By \cref{def:gr_glued_scheme} the diagonal preimage is the overlap
  \(\Delta^{-1}(U^I \times_{\mathbb{Z}} U^J) = U^I \cap U^J = U^I_J = \Spec
  \mathbb{Z}[X^I, 1/P^I_J]\). The restricted diagonal corresponds to the ring map
  \[
    \delta_{I,J} : \mathbb{Z}[X^I] \otimes_{\mathbb{Z}} \mathbb{Z}[X^J]
    \longrightarrow \mathbb{Z}[X^I, 1/P^I_J],
    \qquad
    X^I \otimes 1 \mapsto X^I, \quad 1 \otimes X^J \mapsto (X^I_J)^{-1} X^I,
  \]
  the second component being the comorphism \(\theta_{I,J}\) of
  \cref{def:gr_transition} (a point of \(U^I_J\) and its image under
  \(\theta_{I,J}\) are the two coordinates of a point of the diagonal). We claim
  \(\delta_{I,J}\) is surjective. Its image contains all entries of \(X^I\)
  (from the first factor). From the second factor it contains the entries of
  \((X^I_J)^{-1} X^I\); taking the \(I\)-minor of this matrix, and using
  \(X^I_I = 1_{d\times d}\), the image contains the entries of \((X^I_J)^{-1}\),
  and in particular
  \[
    \delta_{I,J}(1 \otimes P^J_I) = \det\bigl((X^I_J)^{-1}\bigr) = 1/P^I_J .
  \]
  Hence \(1/P^I_J\) lies in the image, so the image is all of
  \(\mathbb{Z}[X^I, 1/P^I_J]\) and \(\delta_{I,J}\) is surjective. Concretely the
  closed subscheme \(\Delta_{I,J} \subseteq U^I \times_{\mathbb{Z}} U^J\) is cut
  out by the entries of the matrix relation
  \[
    X^J_I \, X^I - X^J \;=\; 0,
  \]
  which on the diagonal holds because \(X^J = (X^I_J)^{-1} X^I\) forces
  \(X^J_I = (X^I_J)^{-1} X^I_I = (X^I_J)^{-1}\) and therefore \(X^J_I X^I =
  (X^I_J)^{-1} X^I = X^J\). As each restricted diagonal is a closed immersion,
  \(\Delta\) is a closed immersion and \(\mathrm{Gr}(r,d)\) is separated over
  \(\mathbb{Z}\).
\end{proof}

\section{Properness}
\label{sec:gr_proper}

Following Nitsure \S 1, the structure morphism \(\pi : \mathrm{Gr}(r,d) \to \Spec
\mathbb{Z}\) is proper. We reduce properness to the \emph{existence} half of the
valuative criterion: the Mathlib criterion \texttt{IsProper.of\_valuativeCriterion}
reduces \(\mathrm{IsProper}\,\pi\) to four ingredients --- quasi-compactness,
quasi-separatedness, local finiteness of type, and the valuative criterion
\(\mathrm{Existence} \sqcap \mathrm{Uniqueness}\). The first three and the
uniqueness half are cheap (uniqueness is free from separatedness), so the only
genuine content is \(\mathrm{Existence}\): every commutative square from a
valuation ring admits a diagonal filler. This existence statement is the Nitsure
chart-selection argument, which we decompose into four steps E1--E4.

\subsection*{Cheap ingredients and the reduction to existence}

\begin{lemma}[\(\mathrm{Gr}(r,d)\) is a compact space]\leanok
  \label{lem:gr_compactSpace_scheme}
  \dcref{custom}
  \lean{AlgebraicGeometry.Grassmannian.compactSpace_scheme}
  \uses{def:gr_the_glue_data, def:gr_affine_chart, def:gr_glued_scheme}
  The underlying topological space of the glued Grassmannian scheme
  \(\mathrm{Gr}(r,d)\) (\cref{def:gr_glued_scheme}) is quasi-compact: it carries the
  structure of a \texttt{CompactSpace}.
\end{lemma}
\begin{proof}
  \uses{def:gr_the_glue_data, def:gr_affine_chart}
  \leanok
  The glue datum (\cref{def:gr_the_glue_data}) is indexed by the finite set
  \(\{\, I \subseteq \{1,\dots,r\} : \#(I) = d \,\}\) of size-\(d\) subsets, and
  each chart \(U^I = \Spec \mathbb{Z}[X^I]\) (\cref{def:gr_affine_chart}) is the
  \(\Spec\) of a ring, hence quasi-compact. A scheme glued from finitely many
  quasi-compact charts is quasi-compact, so \(\mathrm{Gr}(r,d)\) is a compact
  space.
\end{proof}

\begin{lemma}[The structure morphism is quasi-compact]\leanok
  \label{lem:gr_quasiCompact_toSpecZ}
  \dcref{custom}
  \lean{AlgebraicGeometry.Grassmannian.quasiCompact_toSpecZ}
  \uses{def:gr_to_specZ, lem:gr_compactSpace_scheme}
  The structure morphism \(\pi : \mathrm{Gr}(r,d) \to \Spec \mathbb{Z}\)
  (\cref{def:gr_to_specZ}) is quasi-compact.
\end{lemma}
\begin{proof}
  \uses{lem:gr_compactSpace_scheme}
  \leanok
  The target \(\Spec \mathbb{Z}\) is affine, so quasi-compactness of \(\pi\) is
  equivalent to quasi-compactness of the total space \(\mathrm{Gr}(r,d)\), which is
  a compact space by \cref{lem:gr_compactSpace_scheme}.
\end{proof}

\begin{lemma}[The structure morphism is locally of finite type]\leanok
  \label{lem:gr_locallyOfFiniteType_toSpecZ}
  \dcref{custom}
  \lean{AlgebraicGeometry.Grassmannian.locallyOfFiniteType_toSpecZ}
  \uses{def:gr_to_specZ, def:gr_the_glue_data, def:gr_affine_chart, lem:gr_ι_toSpecZ}
  The structure morphism \(\pi : \mathrm{Gr}(r,d) \to \Spec \mathbb{Z}\)
  (\cref{def:gr_to_specZ}) is locally of finite type.
\end{lemma}
\begin{proof}
  \uses{def:gr_the_glue_data, lem:gr_ι_toSpecZ}
  \leanok
  Being locally of finite type is local on the source, so it may be checked on the
  chart open cover (\cref{def:gr_the_glue_data}). On the chart \(U^I\) the composite
  \(U^I \to \mathrm{Gr}(r,d) \xrightarrow{\pi} \Spec \mathbb{Z}\) is, by
  \cref{lem:gr_ι_toSpecZ}, the \(\Spec\) of the structure map \(\mathbb{Z} \to R^I
  = \mathbb{Z}[X^I]\). The polynomial ring \(\mathbb{Z}[X^I]\) is a finitely
  generated \(\mathbb{Z}\)-algebra, so this map is of finite type; hence \(\pi\) is
  locally of finite type.
\end{proof}

\begin{lemma}[The structure morphism is quasi-separated]\leanok
  \label{lem:gr_quasiSeparated_toSpecZ}
  \dcref{custom}
  \lean{AlgebraicGeometry.Grassmannian.quasiSeparated_toSpecZ}
  \uses{def:gr_to_specZ, lem:gr_separated_toSpecZ}
  The structure morphism \(\pi : \mathrm{Gr}(r,d) \to \Spec \mathbb{Z}\)
  (\cref{def:gr_to_specZ}) is quasi-separated.
\end{lemma}
\begin{proof}
  \uses{lem:gr_separated_toSpecZ}
  \leanok
  The structure morphism is separated (\cref{lem:gr_separated_toSpecZ}), and every
  separated morphism is quasi-separated.
\end{proof}

\begin{lemma}[Uniqueness half of the valuative criterion]\leanok
  \label{lem:gr_valuativeUniqueness_toSpecZ}
  \dcref{custom}
  \lean{AlgebraicGeometry.Grassmannian.valuativeUniqueness_toSpecZ}
  \uses{def:gr_to_specZ, lem:gr_separated_toSpecZ}
  The structure morphism \(\pi : \mathrm{Gr}(r,d) \to \Spec \mathbb{Z}\)
  (\cref{def:gr_to_specZ}) satisfies the \emph{uniqueness} part of the valuative
  criterion: for a valuation ring \(R\) with fraction field \(K\), any two
  morphisms \(\Spec R \to \mathrm{Gr}(r,d)\) over \(\Spec \mathbb{Z}\) that agree
  on the generic point \(\Spec K\) coincide.
\end{lemma}
\begin{proof}
  \uses{lem:gr_separated_toSpecZ}
  \leanok
  The structure morphism is separated (\cref{lem:gr_separated_toSpecZ}), and the
  uniqueness half of the valuative criterion is free for any separated morphism:
  two lifts agreeing on the schematically dense generic point of \(\Spec R\)
  coincide because the diagonal is a (closed, in particular separated) immersion.
\end{proof}

\begin{lemma}[Properness from the existence half of the valuative criterion]\leanok
  \label{lem:gr_isProper_of_valuativeExistence}
  \dcref{custom}
  \lean{AlgebraicGeometry.Grassmannian.isProper_of_valuativeExistence}
  \uses{lem:gr_compactSpace_scheme, lem:gr_quasiCompact_toSpecZ,
    lem:gr_locallyOfFiniteType_toSpecZ, lem:gr_quasiSeparated_toSpecZ,
    lem:gr_valuativeUniqueness_toSpecZ, def:gr_to_specZ}
  Suppose the structure morphism \(\pi : \mathrm{Gr}(r,d) \to \Spec \mathbb{Z}\)
  (\cref{def:gr_to_specZ}) satisfies the \emph{existence} part of the valuative
  criterion. Then \(\pi\) is proper.
\end{lemma}
\begin{proof}
  \uses{lem:gr_quasiCompact_toSpecZ, lem:gr_locallyOfFiniteType_toSpecZ,
    lem:gr_quasiSeparated_toSpecZ, lem:gr_valuativeUniqueness_toSpecZ}
    \leanok
  The Mathlib criterion \texttt{IsProper.of\_valuativeCriterion} reduces properness
  of \(\pi\) to four hypotheses: \(\pi\) is quasi-compact
  (\cref{lem:gr_quasiCompact_toSpecZ}), quasi-separated
  (\cref{lem:gr_quasiSeparated_toSpecZ}), locally of finite type
  (\cref{lem:gr_locallyOfFiniteType_toSpecZ}), and satisfies the full valuative
  criterion. The full valuative criterion is the conjunction of its existence and
  uniqueness halves; the uniqueness half holds by
  \cref{lem:gr_valuativeUniqueness_toSpecZ}, and the existence half is the
  hypothesis. Assembling the four ingredients gives properness.
\end{proof}

\subsection*{The existence obligation: chart selection by minimal valuation}

The existence half of the valuative criterion is the genuine geometric content,
and it is exactly Nitsure's chart-selection argument. A commutative square over a
valuation ring presents a valuation ring \(R\) with fraction field \(K =
\operatorname{Frac}(R)\), valuation \(v := \mathrm{val}_R : K \to \Gamma\) (written
multiplicatively, so \(R = \{x : v(x) \le 1\}\)), a morphism \(i_1 : \Spec K \to
\mathrm{Gr}(r,d)\), and a morphism \(i_2 : \Spec R \to \Spec \mathbb{Z}\) with
\(i_1 \circ \pi = \Spec(R \to K) \circ i_2\); one must construct a filler
\(\Spec R \to \mathrm{Gr}(r,d)\) making both triangles commute. We carry this out
in four steps. The algebraic core --- the behaviour of the transition hom on a
\(K'\)-minor --- is the following identity over the localised chart ring.

\begin{lemma}[The minor-ratio identity]\leanok
  \label{lem:gr_transitionPreMap_minorDet_mul}
  \dcref{custom}
  \lean{AlgebraicGeometry.Grassmannian.transitionPreMap_minorDet_mul}
  \uses{def:gr_transition_pre, def:gr_minor_det, lem:gr_transitionPreMap_minorDet,
    lem:gr_mul_submatrix_col, lem:gr_universalMinorInv_identities,
    lem:gr_transitionPreMap_minorDet_swap_mul}
  % SOURCE: \cite{nitsure-hilbert-quot}, \S 1, sub-section "Properness" (read from
  % references/nitsure-hilbert-quot-src/nitsure-hilbert-quot.tex, L878-L883).
  % SOURCE QUOTE:
  % "Now consider the homomorphism
  %  $g : \Z[X^J] \to \K$ defined by entries of the matrix formula
  %  $$g(X^J) = f((X^I_J)^{-1}\,X^I)$$
  %  Then $g$ defines the same morphism $\varphi : \spec \K \to \grass(r,d)$,
  %  and moreover all $d\times d$ minors $X^J_K$ satisfy
  %  $\nu(g(P^J_K)) \ge 0$."
  \textit{Source: \cite{nitsure-hilbert-quot}, \S 1.}
  For size-\(d\) subsets \(I, J, K\), the pre-localisation transition hom
  \(\tilde\theta_{I,J}\) (\cref{def:gr_transition_pre}) satisfies, in \(R^I_J =
  \mathbb{Z}[X^I, 1/P^I_J]\), the minor-ratio identity
  \[
    \tilde\theta_{I,J}\bigl(P^J_K\bigr)\cdot P^I_J \;=\; P^I_K
  \]
  (where \(P^I_J, P^I_K\) are read through the structure map \(\mathbb{Z}[X^I] \to
  R^I_J\)). Equivalently \(\tilde\theta_{I,J}(P^J_K) = P^I_K / P^I_J\), the
  determinant of the \(K\)-minor of \((X^I_J)^{-1} X^I\). This generalises
  \cref{lem:gr_transitionPreMap_minorDet_swap_mul} (the case \(K = I\), where
  \(P^I_I = 1\)) and is the algebraic engine of the existence argument: pulled back
  through a \(K\)-point \(f\) it gives \(g(P^J_K) = f(P^I_K)/f(P^I_J)\).
\end{lemma}
\begin{proof}
  \uses{lem:gr_transitionPreMap_minorDet, lem:gr_mul_submatrix_col,
    lem:gr_universalMinorInv_identities}
    \leanok
  By \cref{lem:gr_transitionPreMap_minorDet}, \(\tilde\theta_{I,J}(P^J_K) =
  \det(M_K)\) is the determinant of the \(K\)-minor of the image matrix \(M =
  (X^I_J)^{-1} X^I\). Splitting the \(K\)-column submatrix of the product
  (\cref{lem:gr_mul_submatrix_col}) and using multiplicativity of the determinant,
  \(\det(M_K) = \det\bigl((X^I_J)^{-1}\bigr)\cdot \det(X^I_K) =
  \det(X^I_J)^{-1}\cdot P^I_K\). Multiplying by \(P^I_J = \det(X^I_J)\) and
  cancelling \(\det(X^I_J)^{-1}\det(X^I_J) = 1\)
  (\cref{lem:gr_universalMinorInv_identities}) gives
  \(\tilde\theta_{I,J}(P^J_K)\cdot P^I_J = P^I_K\).
\end{proof}

\begin{lemma}[E1 --- the \(K\)-point factors through a single chart]\leanok
  \label{lem:gr_existence_chart_factorization}
  \dcref{custom}
  \lean{AlgebraicGeometry.Grassmannian.existence_chart_factorization}
  \uses{def:gr_the_glue_data, def:gr_glued_scheme, def:gr_affine_chart,
    lem:gr_chartIncl_isOpenImmersion}
  % SOURCE: \cite{nitsure-hilbert-quot}, \S 1, sub-section "Properness" (read from
  % references/nitsure-hilbert-quot-src/nitsure-hilbert-quot.tex, L868-L870).
  % SOURCE QUOTE:
  % "Let $\Rr$ be a dvr, $\K$ its quotient field, and let
  %  $\varphi : \spec \K \to \grass(r,d)$ be a morphism. This is given by
  %  a ring homomorphism $f : \Z[X^I] \to \K$ for some $I$."
  \textit{Source: \cite{nitsure-hilbert-quot}, \S 1.}
  The morphism \(i_1 : \Spec K \to \mathrm{Gr}(r,d)\) factors through a single
  chart: there exist a size-\(d\) subset \(I\) and a ring homomorphism \(f : R^I =
  \mathbb{Z}[X^I] \to K\) such that
  \[
    i_1 \;=\; \Spec(f) \circ \iota_I : \Spec K \to U^I \to \mathrm{Gr}(r,d),
  \]
  where \(\iota_I : U^I \to \mathrm{Gr}(r,d)\) is the chart immersion of the glue
  datum (\cref{def:gr_the_glue_data}).
\end{lemma}
\begin{proof}
  \uses{def:gr_the_glue_data, def:gr_glued_scheme, lem:gr_chartIncl_isOpenImmersion}
  \leanok
  Since \(K\) is a field, \(\Spec K\) is a one-point space, so its image point \(x
  := i_1(\ast) \in \mathrm{Gr}(r,d)\) is a single point. The chart immersions
  \(\{\iota_I : U^I \to \mathrm{Gr}(r,d)\}\) of the glue datum
  (\cref{def:gr_the_glue_data}) are jointly surjective (their images cover the
  glued scheme \cref{def:gr_glued_scheme}), so \(x\) lies in the range of some
  \(\iota_I\). Each \(\iota_I\) is an open immersion
  (\cref{lem:gr_chartIncl_isOpenImmersion} gives this for the chart inclusions, and
  the glue datum's immersions inherit it). A morphism out of \(\Spec K\) whose
  image point lies in the range of an open immersion factors uniquely through that
  open immersion; hence \(i_1 = j \circ \iota_I\) for a unique \(j : \Spec K \to
  U^I\). As \(U^I = \Spec R^I\) is affine, \(j = \Spec(f)\) for a unique ring map
  \(f : R^I \to K\), giving the claimed factorization.
\end{proof}

\begin{lemma}[E2 --- minimal-valuation chart selection]\leanok
  \label{lem:gr_existence_minimal_valuation}
  \dcref{custom}
  \lean{AlgebraicGeometry.Grassmannian.existence_minimal_valuation}
  \uses{lem:gr_existence_chart_factorization, def:gr_minor_det,
    lem:gr_minorDet_self}
  % SOURCE: \cite{nitsure-hilbert-quot}, \S 1, sub-section "Properness" (read from
  % references/nitsure-hilbert-quot-src/nitsure-hilbert-quot.tex, L870-L876).
  % SOURCE QUOTE:
  % "Having fixed
  %  one such $I$, next choose $J$ such that $\nu(f(P^I_J))$ is minimum,
  %  where $\nu : \K \to \Z\bigcup \{ \infty\}$ denotes
  %  the discrete valuation. As $P^I_I = 1$, note that
  %  $\nu(f(P^I_J)) \le 0$, therefore
  %  $f(P^I_J) \ne 0$ in $\K$ and so the matrix $f(X^I_J)$
  %  lies in $GL_d(\K)$."
  % LEAN SIGNATURE (intended scaffold target). With `v := ValuationRing.valuation
  %   R K`, choose `J : {J : Finset (Fin r) // J.card = d}` maximizing
  %   `v (f (minorDet d r I J))` over the finite index set, and conclude
  %   `f (minorDet d r I J) ≠ 0` (equivalently `v (f (minorDet d r I J)) ≥ 1`).
  \textit{Source: \cite{nitsure-hilbert-quot}, \S 1.}
  Fix \(I\) and \(f : R^I \to K\) as in \cref{lem:gr_existence_chart_factorization},
  and let \(v := \mathrm{val}_R : K \to \Gamma\) be the (multiplicative) valuation
  of \(R\). Over the finite index set \(\{\, J : \#(J) = d \,\}\) choose \(J\)
  \emph{maximising} \(v\bigl(f(P^I_J)\bigr)\) (Nitsure's minimal additive valuation
  \(\nu\) is the maximal multiplicative \(v\)). Then \(I\) itself is a candidate and
  \(P^I_I = 1\) (\cref{lem:gr_minorDet_self}) gives \(v(f(P^I_I)) = v(1) = 1\), so
  the maximum satisfies \(v(f(P^I_J)) \ge 1 > 0\). In particular \(f(P^I_J) \ne 0\),
  so the matrix \(f(X^I_J)\) is invertible over \(K\).
\end{lemma}
\begin{proof}
  \uses{lem:gr_minorDet_self, def:gr_minor_det}
  \leanok
  The index set of size-\(d\) subsets is finite and nonempty (\(I\) is a member),
  so the function \(J \mapsto v(f(P^I_J))\) attains a maximum at some \(J\). Since
  \(P^I_I = 1\) (\cref{lem:gr_minorDet_self}) and \(v(1) = 1\), the value at \(I\)
  is \(1\); hence the maximum is \(\ge 1\), so \(v(f(P^I_J)) \ne 0\) in \(\Gamma\),
  i.e.\ \(f(P^I_J) \ne 0\). A square matrix over the field \(K\) with nonzero
  determinant \(f(P^I_J) = \det f(X^I_J)\) is invertible.
\end{proof}

\begin{lemma}[Determinant of the identity with one column replaced]\leanok
  \label{lem:gr_det_one_updateCol}
  \dcref{custom}
  \lean{AlgebraicGeometry.Grassmannian.det_one_updateCol}
  For a vector \(v : \mathrm{Fin}\,d \to R\) and a row index \(p\), the determinant of the identity
  matrix with its \(p\)-th column replaced by \(v\) is the entry \(v_p\), with \emph{no} sign:
  \(\det\bigl((1).\mathrm{updateCol}\,p\,v\bigr) = v_p\).
\end{lemma}
\begin{proof}

\leanok
  By Cramer's rule \(\det((1).\mathrm{updateCol}\,p\,v) = ((1).\mathrm{cramer}\,v)_p\)
  (\(\mathrm{Matrix.cramer\_apply}\)); and \((1).\mathrm{cramer}\,v = v\) since
  \(\mathrm{1}\cdot\mathrm{cramer}\,v = \det(1)\cdot v = v\) (\(\mathrm{Matrix.mulVec\_cramer}\)).
  Proved directly in Lean.
\end{proof}

\begin{lemma}[Free entry as a signed minor (E3 cofactor sub-step)]\leanok
  \label{lem:gr_free_entry_eq_signed_minor}
  \dcref{custom}
  \lean{AlgebraicGeometry.Grassmannian.exists_minorDet_eq_free_entry}
  \uses{def:gr_universal_matrix, def:gr_minor_det, lem:gr_det_one_updateCol}
  \textit{Source: \cite{nitsure-hilbert-quot}, \S 1 (the cofactor sub-step of the Properness argument).}
  Fix a size-\(d\) subset \(J\), a row index \(p : \mathrm{Fin}\,d\), and a column \(q \notin J\). Put
  \(K' = (J \setminus \{j_p\}) \cup \{q\}\), again of size \(d\). Then there is an order-iso realisation of
  \(K'\) under which the \(K'\)-minor of the universal matrix equals, up to sign, the free entry:
  \[
    \mathrm{minorDet}\,d\,r\,J\,K' \;=\; \pm\, X_{(p,\,q)} .
  \]
\end{lemma}
\begin{proof}
  \uses{def:gr_universal_matrix, def:gr_minor_det, lem:gr_det_one_updateCol}
  \leanok
  Reorder the columns of \(K'\) so that column \(q\) sits in the \(p\)-th slot and the remaining columns
  match \(J\) in increasing order; call \(\sigma\) the resulting permutation. Because the \(J\)-minor
  \(X^J_J\) is the identity, the so-reordered submatrix is the identity with its \(p\)-th column overwritten
  by the single free column carrying \(X_{(p,q)}\). The determinant of an identity matrix with one column
  replaced by a vector \(v\) is exactly \(v_p\), with \emph{no} sign --- via Cramer's rule,
  \(\det\bigl((1).\mathrm{updateCol}\,p\,v\bigr) = \bigl((1).\mathrm{cramer}\,v\bigr)_p
  = (\mathrm{1}\cdot v)_p = v_p\). Restoring the increasing column order multiplies by \(\mathrm{sign}(\sigma)
  = \pm 1\) (\(\mathrm{det\_permute}'\)), so \(\mathrm{minorDet}\,d\,r\,J\,K' = \mathrm{sign}(\sigma)\cdot
  X_{(p,q)} = \pm X_{(p,q)}\).
\end{proof}

\begin{lemma}[E3 --- entries land in \(R\); \(g\) factors through the valuation ring]\leanok
  \label{lem:gr_existence_factor_through_valuation_ring}
  \dcref{custom}
  \lean{AlgebraicGeometry.Grassmannian.existence_factor_through_valuationRing}
  \uses{lem:gr_transitionPreMap_minorDet_mul, lem:gr_existence_minimal_valuation,
    lem:gr_transitionPreMap_minorDet, lem:gr_isUnit_incl_transitionPreMap_cross,
    lem:gr_free_entry_eq_signed_minor,
    lem:mathlib_away_lift, def:gr_transition_pre, def:gr_minor_det}
  % SOURCE: \cite{nitsure-hilbert-quot}, \S 1, sub-section "Properness" (read from
  % references/nitsure-hilbert-quot-src/nitsure-hilbert-quot.tex, L881-L887).
  % SOURCE QUOTE:
  % "Then $g$ defines the same morphism $\varphi : \spec \K \to \grass(r,d)$,
  %  and moreover all $d\times d$ minors $X^J_K$ satisfy
  %  $\nu(g(P^J_K)) \ge 0$.
  %  As the minor $X^J_J$ is identity, it follows from the above that in fact
  %  $\nu(g(x^J_{p,q}))\ge 0$ for all entries of $X^J$.
  %  Therefore, the map $g : \Z[X^J] \to \K$
  %  factors uniquely via $\Rr \subset \K$."
  % LEAN SIGNATURE (intended scaffold target). Set
  %   `g := f' ∘ transitionPreMap d r I J`, where `f' : R^I_J → K` is the
  %   `IsLocalization.Away.lift` of `f` along `minorDet d r I J` (a unit by E2).
  %   Prove `∀ x, g x ∈ (algebraMap R K).range`, hence factor `g = (algebraMap R K)
  %   ∘ g'` for a unique ring hom `g' : R^J → R`.
  \textit{Source: \cite{nitsure-hilbert-quot}, \S 1.}
  Define \(g : R^J = \mathbb{Z}[X^J] \to K\) by the matrix formula \(g(X^J) =
  f\bigl((X^I_J)^{-1} X^I\bigr)\); concretely \(g = f' \circ \tilde\theta_{I,J}\),
  where \(f' : R^I_J \to K\) is the extension of \(f\) to the away-localisation
  \(R^I_J = \mathbb{Z}[X^I, 1/P^I_J]\) (well-defined by the universal property
  \cref{lem:mathlib_away_lift}, since \(f(P^I_J)\) is a unit of \(K\) by
  \cref{lem:gr_existence_minimal_valuation}). Then every value \(g(P^J_K)\) lies in
  \(R\), and consequently \(g\) factors uniquely as \(g = (R \hookrightarrow K)
  \circ g'\) for a ring homomorphism \(g' : R^J \to R\). Moreover \(g\) defines the
  same \(K\)-point as \(f\).
\end{lemma}
\begin{proof}
  \uses{lem:gr_transitionPreMap_minorDet_mul, lem:gr_existence_minimal_valuation,
    lem:gr_transitionPreMap_minorDet, lem:gr_isUnit_incl_transitionPreMap_cross}
    \leanok
  Applying \(f'\) to the minor-ratio identity
  (\cref{lem:gr_transitionPreMap_minorDet_mul}), and using that \(f'\) extends
  \(f\), gives for every size-\(d\) subset \(K'\)
  \[
    g\bigl(P^J_{K'}\bigr)\cdot f(P^I_J) \;=\; f(P^I_{K'}),
    \qquad\text{so}\qquad
    g\bigl(P^J_{K'}\bigr) \;=\; \frac{f(P^I_{K'})}{f(P^I_J)} .
  \]
  By the maximality in \cref{lem:gr_existence_minimal_valuation},
  \(v(f(P^I_{K'})) \le v(f(P^I_J))\), hence \(v\bigl(g(P^J_{K'})\bigr) =
  v(f(P^I_{K'}))/v(f(P^I_J)) \le 1\), i.e.\ \(g(P^J_{K'}) \in R\). (That
  \(f(P^I_J)\) is a unit and the cross minors map to units is recorded by
  \cref{lem:gr_isUnit_incl_transitionPreMap_cross}, via the determinant identity of
  \cref{lem:gr_transitionPreMap_minorDet}.) Now the \(J\)-minor \(X^J_J\) of \(X^J\)
  is the identity, so cofactor expansion expresses each free entry \(x^J_{p,q}\)
  (\(q \notin J\)) as, up to sign, the minor \(P^J_{K'}\) with \(K' = (J \setminus
  \{j_p\}) \cup \{q\}\) --- the determinant obtained by replacing the \(p\)-th
  identity column of \(X^J_J\) by column \(q\). Hence \(g(x^J_{p,q}) = \pm
  g(P^J_{K'}) \in R\), while the entries with \(q \in J\) are \(0\) or \(1\),
  already in \(R\). (This entry-as-minor sign bookkeeping is the one sub-step with
  no pre-existing matrix-algebra scaffold; it requires expanding the determinant of
  a column-substituted identity matrix.) Since \(g\) sends every generator of
  \(\mathbb{Z}[X^J]\) into \(R\), it factors uniquely through the subring \(R
  \subseteq K\) as \(g = (R \hookrightarrow K) \circ g'\) with \(g' : R^J \to R\).
\end{proof}


\begin{lemma}[E4 top-triangle $K$-point identity]\leanok
  \label{lem:gr_existence_chart_kpoint_eq}
  \dcref{custom}
  \lean{AlgebraicGeometry.Grassmannian.existence_chart_kpoint_eq}
  \uses{lem:gr_chartTransition_comp_chartIncl, def:gr_the_glue_data,
    lem:gr_transitionPreMap_minorDet}
  \textit{Source: \cite{nitsure-hilbert-quot}, \S 1.}
  For charts \(I, J\) of \(\mathrm{Gr}(r,d)\) and the comparison ring map
  \(g := f' \circ \tilde\theta_{I,J} : R^I \to R\) (with \(f'\) the localized lift of
  \(f : R^I \to R\) along the away-localization and \(\tilde\theta_{I,J}\) the chart-transition
  comorphism), the two chart legs present the same \(K\)-point:
  \[
    \Spec(g) \circ \iota_J \;=\; \Spec(f) \circ \iota_I .
  \]
\end{lemma}
\begin{proof}
  \uses{lem:gr_chartTransition_comp_chartIncl, def:gr_the_glue_data}
  \leanok
  Expand \(\Spec(g) = \Spec(\tilde\theta_{I,J}) \circ \Spec(f')\) and rewrite
  \(\Spec(\tilde\theta_{I,J}) = t_{I,J} \circ \iota_{J,I}\) via
  \cref{lem:gr_chartTransition_comp_chartIncl}; the glue cocycle identity
  \(t \circ f \circ \iota = f \circ \iota\) of the glue data (\cref{def:gr_the_glue_data}) together with
  the localized factorization \(f' \) followed by the away map \(= f\) collapses the composite to
  \(\Spec(f) \circ \iota_I\).
\end{proof}

\begin{lemma}[E4 --- the filler and its two triangles]\leanok
  \label{lem:gr_existence_lift}
  \dcref{custom}
  \lean{AlgebraicGeometry.Grassmannian.existence_lift}
  \uses{lem:gr_existence_chart_factorization,
    lem:gr_existence_factor_through_valuation_ring, def:gr_the_glue_data,
    def:gr_to_specZ, lem:gr_ι_toSpecZ, lem:gr_existence_chart_kpoint_eq}
  % SOURCE: \cite{nitsure-hilbert-quot}, \S 1, sub-section "Properness" (read from
  % references/nitsure-hilbert-quot-src/nitsure-hilbert-quot.tex, L887-L889).
  % SOURCE QUOTE:
  % "The resulting morphism of
  %  schemes $\spec \Rr \to U^J \hra \grass(r,d)$ prolongs
  %  $\varphi : \spec \K \to \grass(r,d)$."
  % LEAN SIGNATURE (intended scaffold target). The filler is
  %   `Spec.map (CommRingCat.ofHom g') ≫ (theGlueData d r).ι J`; assemble a
  %   `CommSq.LiftStruct` (top triangle from E1 + `g = algebraMap ∘ g'` + ι_toSpecZ,
  %   bottom triangle from `specZIsTerminal.hom_ext`) and conclude `HasLift`.
  \textit{Source: \cite{nitsure-hilbert-quot}, \S 1.}
  The morphism
  \[
    \ell \;:=\; \Spec(g') \circ \iota_J : \Spec R \to U^J \to \mathrm{Gr}(r,d)
  \]
  (with \(g' : R^J \to R\) from
  \cref{lem:gr_existence_factor_through_valuation_ring} and \(\iota_J\) the chart
  immersion of \cref{def:gr_the_glue_data}) is a diagonal filler of the valuative
  square: it restricts to \(i_1\) over the generic point \(\Spec K\), and it lies
  over \(\Spec \mathbb{Z}\). Hence the valuative square has a lift.
\end{lemma}
\begin{proof}
  \uses{lem:gr_existence_chart_factorization,
    lem:gr_existence_factor_through_valuation_ring, lem:gr_ι_toSpecZ,
    def:gr_to_specZ}
    \leanok
  \emph{Top triangle.} Restricting \(\ell\) along \(\Spec(R \to K)\) yields
  \(\Spec(g' \text{ followed by } R \hookrightarrow K) \circ \iota_J = \Spec(g)
  \circ \iota_J\), using \(g = (R \hookrightarrow K)\circ g'\)
  (\cref{lem:gr_existence_factor_through_valuation_ring}). Since \(g = f' \circ
  \tilde\theta_{I,J}\) presents the same \(K\)-point as \(f\), and \(i_1 = \Spec(f)
  \circ \iota_I\) (\cref{lem:gr_existence_chart_factorization}), this composite
  equals \(i_1\). \emph{Bottom triangle.} Both \(\ell \circ \pi\) and the given
  \(i_2\) are morphisms \(\Spec R \to \Spec \mathbb{Z}\) into the terminal object
  (\cref{def:gr_to_specZ}; \cref{lem:gr_ι_toSpecZ} identifies the chart leg with the
  affine structure map), hence coincide by uniqueness of the terminal map. The two
  triangles assemble into a lift structure for the valuative square, so a lift
  exists.
\end{proof}

\begin{lemma}[Existence half of the valuative criterion]\leanok
  \label{lem:gr_valuativeExistence_toSpecZ}
  \dcref{custom}
  \lean{AlgebraicGeometry.Grassmannian.valuativeExistence_toSpecZ}
  \uses{def:gr_to_specZ, lem:gr_existence_chart_factorization,
    lem:gr_existence_minimal_valuation,
    lem:gr_existence_factor_through_valuation_ring, lem:gr_existence_lift}
  % SOURCE: \cite{nitsure-hilbert-quot}, \S 1, sub-section "Properness" (read from
  % references/nitsure-hilbert-quot-src/nitsure-hilbert-quot.tex, L866-L889).
  % SOURCE QUOTE:
  % "We now show that $\pi : \grass(r,d) \to \spec \Z$ is proper. It is enough to
  %  verify the valuative criterion of properness for discrete valuation rings.
  %  ... The resulting morphism of
  %  schemes $\spec \Rr \to U^J \hra \grass(r,d)$ prolongs
  %  $\varphi : \spec \K \to \grass(r,d)$."
  \textit{Source: \cite{nitsure-hilbert-quot}, \S 1.}
  The structure morphism \(\pi : \mathrm{Gr}(r,d) \to \Spec \mathbb{Z}\)
  (\cref{def:gr_to_specZ}) satisfies the \emph{existence} part of the valuative
  criterion: every commutative square from a valuation ring \(R\) (with fraction
  field \(K\)) admits a diagonal filler \(\Spec R \to \mathrm{Gr}(r,d)\).
\end{lemma}
\begin{proof}
  \uses{lem:gr_existence_chart_factorization, lem:gr_existence_minimal_valuation,
    lem:gr_existence_factor_through_valuation_ring, lem:gr_existence_lift}
    \leanok
  Given a valuative square, factor the generic-point morphism through a chart
  (\cref{lem:gr_existence_chart_factorization}), select \(J\) by maximal valuation
  (\cref{lem:gr_existence_minimal_valuation}), factor the transported ring map
  \(g\) through the valuation ring (\cref{lem:gr_existence_factor_through_valuation_ring}),
  and take the filler \(\Spec(g') \circ \iota_J\)
  (\cref{lem:gr_existence_lift}). This produces the required diagonal filler, so the
  existence half holds.
\end{proof}

\subsection*{Properness}

\begin{lemma}[\(\mathrm{Gr}(r,d)\) is proper over \(\mathbb{Z}\)]\leanok
  \label{lem:gr_proper}
  \dcref{custom}
  \lean{AlgebraicGeometry.Grassmannian.isProper}
  \uses{def:gr_to_specZ, lem:gr_isProper_of_valuativeExistence,
    lem:gr_valuativeExistence_toSpecZ}
  % SOURCE: \cite{nitsure-hilbert-quot}, \S 1, sub-section "Properness" (read from
  % references/nitsure-hilbert-quot-src/nitsure-hilbert-quot.tex, L865-L891).
  % SOURCE QUOTE:
  % "We now show that $\pi : \grass(r,d) \to \spec \Z$ is proper. It is enough to
  %  verify the valuative criterion of properness for discrete valuation rings.
  %  Let $\Rr$ be a dvr, $\K$ its quotient field, and let
  %  $\varphi : \spec \K \to \grass(r,d)$ be a morphism. This is given by
  %  a ring homomorphism $f : \Z[X^I] \to \K$ for some $I$. Having fixed
  %  one such $I$, next choose $J$ such that $\nu(f(P^I_J))$ is minimum,
  %  where $\nu : \K \to \Z\bigcup \{ \infty\}$ denotes
  %  the discrete valuation. As $P^I_I = 1$, note that
  %  $\nu(f(P^I_J)) \le 0$, therefore
  %  $f(P^I_J) \ne 0$ in $\K$ and so the matrix $f(X^I_J)$
  %  lies in $GL_d(\K)$.
  %
  %  Now consider the homomorphism
  %  $g : \Z[X^J] \to \K$ defined by entries of the matrix formula
  %  $$g(X^J) = f((X^I_J)^{-1}\,X^I)$$
  %  Then $g$ defines the same morphism $\varphi : \spec \K \to \grass(r,d)$,
  %  and moreover all $d\times d$ minors $X^J_K$ satisfy
  %  $\nu(g(P^J_K)) \ge 0$.
  %  As the minor $X^J_J$ is identity, it follows from the above that in fact
  %  $\nu(g(x^J_{p,q}))\ge 0$ for all entries of $X^J$.
  %  Therefore, the map $g : \Z[X^J] \to \K$
  %  factors uniquely via $\Rr \subset \K$. The resulting morphism of
  %  schemes $\spec \Rr \to U^J \hra \grass(r,d)$ prolongs
  %  $\varphi : \spec \K \to \grass(r,d)$. We have already checked
  %  separatedness of $\grass(r,d)$, so now we see that
  %  $\grass(r,d) \to \spec \Z$ is proper."
  \textit{Source: \cite{nitsure-hilbert-quot}, \S 1.}
  The structure morphism \(\pi : \mathrm{Gr}(r,d) \to \Spec \mathbb{Z}\)
  (\cref{def:gr_glued_scheme}) is proper.
\end{lemma}

% SOURCE QUOTE PROOF: "We now show that $\pi : \grass(r,d) \to \spec \Z$ is
% proper. It is enough to verify the valuative criterion of properness for
% discrete valuation rings. Let $\Rr$ be a dvr, $\K$ its quotient field, and let
% $\varphi : \spec \K \to \grass(r,d)$ be a morphism. This is given by a ring
% homomorphism $f : \Z[X^I] \to \K$ for some $I$. Having fixed one such $I$, next
% choose $J$ such that $\nu(f(P^I_J))$ is minimum, where
% $\nu : \K \to \Z\bigcup \{ \infty\}$ denotes the discrete valuation. As
% $P^I_I = 1$, note that $\nu(f(P^I_J)) \le 0$, therefore $f(P^I_J) \ne 0$ in $\K$
% and so the matrix $f(X^I_J)$ lies in $GL_d(\K)$. Now consider the homomorphism
% $g : \Z[X^J] \to \K$ defined by entries of the matrix formula
% $$g(X^J) = f((X^I_J)^{-1}\,X^I)$$ Then $g$ defines the same morphism
% $\varphi : \spec \K \to \grass(r,d)$, and moreover all $d\times d$ minors
% $X^J_K$ satisfy $\nu(g(P^J_K)) \ge 0$. As the minor $X^J_J$ is identity, it
% follows from the above that in fact $\nu(g(x^J_{p,q}))\ge 0$ for all entries of
% $X^J$. Therefore, the map $g : \Z[X^J] \to \K$ factors uniquely via
% $\Rr \subset \K$. The resulting morphism of schemes
% $\spec \Rr \to U^J \hra \grass(r,d)$ prolongs
% $\varphi : \spec \K \to \grass(r,d)$. We have already checked separatedness of
% $\grass(r,d)$, so now we see that $\grass(r,d) \to \spec \Z$ is proper."
% (references/nitsure-hilbert-quot-src/nitsure-hilbert-quot.tex, L865-L891.)
\begin{proof}
  \uses{lem:gr_isProper_of_valuativeExistence, lem:gr_valuativeExistence_toSpecZ}
  \leanok
  Nitsure's argument is precisely the valuative criterion, organised here so that
  all the topological/finiteness bookkeeping is separated from the one geometric
  step. The reduction \cref{lem:gr_isProper_of_valuativeExistence} packages
  quasi-compactness (\cref{lem:gr_quasiCompact_toSpecZ}), quasi-separatedness
  (\cref{lem:gr_quasiSeparated_toSpecZ}), local finiteness of type
  (\cref{lem:gr_locallyOfFiniteType_toSpecZ}) and the uniqueness half
  (\cref{lem:gr_valuativeUniqueness_toSpecZ}, free from separatedness) into the
  hypothesis that \(\pi\) satisfies only the \emph{existence} half of the valuative
  criterion. That existence half is established in
  \cref{lem:gr_valuativeExistence_toSpecZ} by the chart-selection argument of steps
  E1--E4: the generic-point morphism lands in a chart \(U^I\) via a ring map \(f\)
  (\cref{lem:gr_existence_chart_factorization}); one selects \(J\) with
  \(\nu(f(P^I_J))\) minimal --- equivalently \(v(f(P^I_J))\) maximal ---
  (\cref{lem:gr_existence_minimal_valuation}), so \(f(X^I_J) \in \mathrm{GL}_d(K)\);
  the transported map \(g(X^J) = f((X^I_J)^{-1} X^I)\) then sends every generator
  into the valuation ring and factors as \(g = (R \hookrightarrow K)\circ g'\)
  (\cref{lem:gr_existence_factor_through_valuation_ring}); and \(\Spec(g') \circ
  \iota_J\) is the prolonging filler (\cref{lem:gr_existence_lift}). Feeding this
  existence statement into \cref{lem:gr_isProper_of_valuativeExistence} yields that
  \(\pi : \mathrm{Gr}(r,d) \to \Spec \mathbb{Z}\) is proper.
\end{proof}

\section{Out of scope}
\label{sec:gr_out_of_scope}

This chapter constructs the Grassmannian \emph{scheme} \(\mathrm{Gr}(r,d)\) over
the absolute base \(\mathbb{Z}\) and establishes its smoothness, separatedness
and properness. The following constructions are deferred to later chapters and
are not treated here:
\begin{itemize}
  \item \textbf{The tautological/universal quotient} \(u : \mathcal{O}^{\oplus
    r}_{\mathrm{Gr}(r,d)} \twoheadrightarrow \mathcal{U}\) and the determinant
    line bundle \(\det \mathcal{U}\) with its Pl\"ucker embedding.
  \item \textbf{The functor-of-points representability} of \(\mathrm{Gr}(r,d)\)
    -- that it represents \(\mathrm{Grass}(r,d) =
    \Quot^{d,\mathcal{O}_{\mathbb{Z}}}_{\oplus^r \mathcal{O}_{\mathbb{Z}} /
    \mathbb{Z} / \mathbb{Z}}\); this is blueprinted as
    \cref{thm:grassmannian_representable} in
    \cref{chap:Picard_QuotScheme}.
  \item \textbf{The relative version over a general base \(S\)} with a locally
    free sheaf \(V\) of rank \(r\): the scheme \(\mathrm{Gr}_S(V,d)\) obtained by
    base change and gluing over a trivialising cover of \(V\). The
    absolute-over-\(\mathbb{Z}\) construction of this chapter is the input to
    that base change.
\end{itemize}
